% concatenated from Endo_SL2_2025_Revised_v3.tex
\documentclass[11pt]{amsart}

%\makeatletter
%t\@wraptoccontribs\wraptoccontribs
%\makeatother

\usepackage{latexsym}
\usepackage[utf8]{inputenc}
\usepackage{amsmath}
\usepackage{amssymb}
\usepackage{amsthm}
\usepackage{amsfonts}
\usepackage{mathrsfs}

\usepackage{hyperref} 
%\usepackage{enumerate}
\usepackage{graphicx}

%\usepackage{showlabels}
\usepackage[all]{xy}  
\usepackage{mathtools}
%\usepackage{refcheck}
\usepackage{color}
\usepackage{comment}
\usepackage{appendix}
\usepackage{tikz}
\usepackage{tikz-cd}
\usepackage[english]{babel}
\usepackage{amscd,url,geometry,pdflscape}
\usepackage{comment}
\usepackage{mathalfa}
%\usepackage{verbatim,url}

% -- Jessica added:
\usepackage{marginnote}
\newcommand{\mar}[1]{\marginpar{\raggedright\tiny #1}} %tiny margin remarks

%%% ---
 

\setlength{\unitlength}{1cm}
\setlength{\topmargin}{0cm}
\setlength{\textheight}{22cm}
\setlength{\oddsidemargin}{1cm}
\setlength{\textwidth}{14cm}
\setlength{\voffset}{-1cm}

\numberwithin{equation}{section}

\newtheorem{thm}[equation]{Theorem}
\newtheorem{cor}[equation]{Corollary}
\newtheorem{lem}[equation]{Lemma}
\newtheorem{prop}[equation]{Proposition}
\newtheorem{defn}[equation]{Definition}
\newtheorem{conj}[equation]{Conjecture}
\newtheorem{cond}[equation]{Condition}
\newtheorem{rem}[equation]{Remark}
\newtheorem{thmintro}{Theorem}
\newtheorem{ex}[equation]{Example}
\def\cay{\mathrm{cay}}
\def\kF{\mathfrak{f}}
\def\Fq{{\mathbb{F}_q}}
\def\Fqbar{\overline{\mathbb{F}}_q}
\def\FF{{\mathbb{F}}}
\def\cInd{\mathrm{c\!\!-\!\!Ind}}
\def\rss{\mathrm{rss}}
\def\LLC{\mathrm{LLC}}
\def\HII{\mathrm{HII}}
\def\Res{\mathrm{Res}}
\def\bd{\mathrm{bd}}
\def\Nrd{\mathrm{Nrd}}
\def\g{\mathbf{g}}
\def\ppi{\mathbf{\pi}}
\def\Irr{\mathrm{Irr}}
\def\sgn{\mathrm{sgn}}
\def\bG{\mathcal{G}}
\def\cO{\mathcal{O}}
\def\I{\mathfrak{I}}
\def\cB{\mathcal{B}}
\def\fB{\mathfrak{B}}
\def\fL{\mathfrak{L}}
\def\fR{\mathfrak{R}}
\def\cI{\mathcal{I}}
\def\cP{\mathcal{P}}
\def\cH{\mathcal{H}}
\def\cS{\mathcal{S}}
\def\cT{\mathcal{T}}
\def\cU{\mathcal{U}}
\def\cW{\mathcal{W}}
\def\fD{\mathfrak{D}}
\def\JL{\mathfrak{JL}}
\def\fg{\mathfrak{g}}
\def\fh{\mathfrak{h}}
\def\fs{\mathfrak{s}}
\def\fo{\mathfrak{o}}
\def\fp{\mathfrak{p}}
\def\ft{\mathfrak{t}}
\def\fT{\mathfrak{T}}
\def\fR{\mathfrak{R}}
\def\fe{\mathfrak{e}}
\def\fw{\mathfrak{w}}
\def\fsl{\mathfrak{sl}}
\def\fso{\mathfrak{so}}
\def\adm{\mathrm{adm}}
\def\Ad{\mathrm{Ad}}
\def\ad{\mathrm{ad}}
\def\gen{\mathrm{g}}
\def\sr{\mathrm{sr}}
\def\sc{\mathrm{\sco}}
\def\card{\mathrm{card}}
\def\pr{\mathrm{pr}}
\def\Pl{\mathrm{Pl}}
\def\rH{\mathrm{H}}
\def\rO{\mathrm{O}}
\def\rU{\mathrm{U}}
\def\rZ{\mathrm{Z}}
\def\Ru{{\mathrm{R}_{\mathrm u}}}
%\def\char{\mathrm{char}}
\def\cl{\mathrm{cl}}
\def\Cent{\mathrm{Cent}}
\def\Deg{\mathrm{Deg}}
\def\dim{\mathrm{dim}}
\def\Gal{\mathrm{Gal}}
\def\Hom{\mathrm{Hom}}
\def\H{\mathrm{H}}
\def\im{\mathrm{im}}
\def\nr{\mathrm{nr}}
\def\Ind{\mathrm{Ind}}
\def\rI{\mathrm{I}}
\def\ind{\mathrm{ind}}
\def\KLR{\mathrm{KLR}}
\def\Mat{\mathrm{M}}
\def\Mod{\mathrm{Mod}}
\def\Nor{\mathrm{N}}
\def\cusp{\mathrm{cusp}}
\def\dep{\mathrm{dep}}
\def\span{\mathrm{span}}
\def\temp{\mathrm{temp}}
\def\St{\mathrm{St}}
\def\sep{\mathrm{s}}
\def\uur{\mathrm{uur}}
\def\ur{\mathrm{ur}}
\def\ab{\mathrm{ab}}
\def\enh{\mathrm{e}}
\def\fdeg{\mathrm{fdeg}}
\def\val{\mathrm{v}}
\def\reg{\mathrm{reg}}
\def\JL{\mathrm{JL}}
\def\LQ{\mathrm{LQ}}
\def\SO{\mathrm{SO}}
\def\SU{\mathrm{SU}}
\newcommand{\isom}{\xrightarrow{\;\sim\;}}
\def\PGL{\mathrm{PGL}}
\def\PSU{\mathrm{PSU}}
\def\GL{\mathrm{GL}}
\def\Aut{\mathrm{Aut}}
\def\Nrd{\mathrm{Nrd}}
\def\F{\mathbb{F}}
\def\bW{\mathbf{W}}
\def\R{\mathbb{R}}
\def\C{\mathbb{C}}
\def\N{\mathbb{N}}
\def\Z{\mathbb{Z}}
\def\Q{\mathbb{Q}}
\def\Qp{\mathbb{Q}_p}
\def\SL{\mathrm{SL}}
\def\Sh{\mathrm{Sh}}
\def\Int{\mathrm{Int}}
\def\Out{\mathrm{Out}}
\def\Id{\mathrm{Id}}
\def\diag{\mathrm{diag}}
\def\PSL{\mathrm{PSL}}
\def\im{\mathrm{im}}
\def\mini{\mathrm{mini}}
\def\max{\mathrm{max}}
\def\triv{\mathrm{triv}}
\def\res{\mathrm{res}}
\def\rec{\mathrm{rec}}
\def\sep{\mathrm{sep}}
\def\Tr{\mathrm{Tr}}
\def\Map{\mathrm{Map}}
\def\Prim{\mathbf{Prim}}
\def\I{\mathbf{I}}
\def\i{\mathbf{i}}
\def\j{\mathbf{j}}
\def\k{\mathbf{k}}
\def\h{\mathbf{h}}
\def\M{\mathbf{M}}
\def\n{\mathbf{n}}
\def\bS{\mathbf{S}}
\def\U{\mathbf{U}}
\def\q{/\!/}
\def\d{{\rm d}}
\def\sep{{\mathrm{s}}}
\def\gphi{{\varphi}}
\newcommand{\dv}{d_{\varpi}}
\newcommand{\mf}{\mathfrak}
\newcommand{\mc}{\mathcal}
\def\barF{{F_\sep}}
\def\bark{{k_\sep}}
\def\cA{\mathcal{A}}
\def\cC{\mathcal{C}}
\def\cE{\mathcal{E}}
\def\cF{\mathcal{F}}
\def\cX{\mathcal{X}}
\def\cK{\mathcal{K}}
\def\cR{\mathcal{R}}
\def\cZ{\mathcal{Z}}
\def\tG{\widetilde{{G}}}
\def\bbG{\mathbb{G}}
\def\M{\mathbf{M}}
\def\bHom{\mathbf{Hom}}
\def\bG{\mathbf{G}}
\def\tbG{\widetilde{{\bG}}}
\def\bI{\mathbf{I}}
\def\bB{\mathbf{B}}
\def\bM{\mathbf{M}}
\def\bH{\mathbf{H}}
\def\bT{\mathbf{T}}
\def\bbT{\mathbb{T}}
\def\bZ{\mathbf{Z}}
\def\nr{\mathrm{nr}}
\def\As{\mathrm{As}}
\def\rmA{\mathrm{A}}
\def\red{\mathrm{red}}
\def\Stab{\mathrm{Stab}}
\def\Gm{{\mathbb{G}_{\mathrm m}}}
\def\W{\mathbf{W}}
\def\fa{\mathfrak{a}}
\def\fo{\mathfrak{o}}
\def\fp{\mathfrak{p}}
\def\fR{\mathfrak{R}}
\def\fS{\mathfrak{S}}
\def\fZ{\mathfrak{Z}}
\def\bdd{\mathrm{bdd}}
\def\canon{{\mathrm{can}}}
\def\depth{{\mathrm{depth}}}
\def\en{\mathrm{e}}
\def\Int{{\mathrm{Int}}}
\def\id{{\mathrm{id}}}
\def\LLC{\mathrm{{LLC}}}
\def\rM{{\mathrm{M}}}
\def\Nor{{\mathrm{N}}}
\def\ii{{\mathrm{i}}}
\def\Ob{{\mathrm{Ob}}}
\def\rig{{\mathrm{\rig}}}
\def\Sch{\mathrm{Sch}}
\def\fppf{{\mathrm{fppf}}}
\def\fpqc{{\mathrm{fpqc}}}
\def\spl{{\mathrm{spl}}}
\def\sign{\mathrm{sign}}
\def\sw{\mathrm{sw}}
\def\temp{{\mathrm{t}}}
\def\tr{{\mathrm{tr}}}
\def\Lift{\mathrm{Lift}}
\def\Trans{\mathrm{Trans}}
\def\trace{\mathrm{trace}}
\def\bas{{\mathrm{bas}}}
\def\pr{\mathrm{pr}}
\def\ad{{\mathrm{ad}}}
\def\der{{\mathrm{der}}}
\def\rig{{\mathrm{rig}}}
\def\rN{{\mathrm{N}}}
\def\rS{{\mathrm{S}}}
\def\sco{{\mathrm{sc}}}
\newcommand{\sym}{\mathrm{sym}}
\newcommand{\matje}[4]{\left(\begin{smallmatrix} #1 & #2 \\ 
#3 & #4 \end{smallmatrix}\right)}

\def\tB{\widetilde{B}}
\def\tT{\widetilde{T}}
\def\ts{\tilde{s}}
\def\teta{\tilde{\eta}}
\def\TGur{\mathbf T^{\Gur}}
\def\TMur{\mathbf T^{\Mur}}
\def\Mur{\M}
\def\Gur{\G}
\def\SGur{\mathbf S^{\Gur}}
\def\SMur{\mathbf S^{\Mur}}
\def\UGur{\mathbf U^{\Gur}}
\def\UMur{\mathbf U^{\Mur}}
\def\xMur{x^{\Mur}}
\def\tM{\widetilde{M}}




\newcommand{\wt}{\widetilde}
\newcommand{\ssc}{\mathscr{C}}
\newcommand{\sB}{\mathscr{B}}
\newcommand{\sA}{\mathscr{A}}
\newcommand{\sG}{\mathscr{G}}
\newcommand{\sH}{\mathscr{H}}
\newcommand{\sM}{\mathscr{M}}
\newcommand{\sR}{\mathscr{R}}
\newcommand{\sS}{\mathscr{S}}
\newcommand{\sZ}{\mathscr{Z}}
\def\bij{{\mathscr{L}}}
\newcommand{\bR}{\mathbb{R}}
\def\a{\mathbf{a}}
\def\I{\mathbf{I}}
\def\i{\mathbf{i}}
\def\j{\mathbf{j}}
\def\k{\mathbf{k}}
\def\h{\mathbf{h}}
\def\a{\mathbf{a}}
\def\I{\mathbf{I}}
\def\i{\mathbf{i}}
\def\j{\mathbf{j}}
\def\k{\mathbf{k}}
\def\h{\mathbf{h}}
\def\g{\mathbf{g}}
\def\u{\mathbf{u}}
\DeclareSymbolFont{bbold}{U}{bbold}{m}{n}
\DeclareSymbolFontAlphabet{\mathbbold}{bbold}
\newcommand{\bbbG}{{\mathbbold{G}}}
\newcommand{\bbbM}{{\mathbbold{M}}}
\newcommand{\bbbS}{{\mathbbold{S}}}
\DeclareMathOperator{\Spec}{Spec}

\newcommand{\ra}{\rightarrow}

%Jessica added:
\newcommand{\coB}{\textcolor{blue}}
\newcommand{\coR}{\textcolor{red}}
\def\T{\mathbf{T}}
\newcommand{\<}{\left\langle}
\renewcommand{\>}{\right\rangle}

\setlength{\parindent}{0pt} % to have no indent for a new paragraph
\setlength{\parskip}{1ex plus 0.5ex minus 0.2ex}



\begin{document}

\title[Endoscopy]{Comparison of the Sally-Shalika character formulas with the endoscopic character identities for  $\SL_2$} 




\author[A.-M. Aubert]{Anne-Marie Aubert}
\address{Sorbonne Universit\'e and Universit\'e Paris Cit\'e, CNRS, IMJ-PRG, 
F-75005 Paris, France}
\email{anne-marie.aubert@imj-prg.fr}
\author[R. Plymen]{Roger Plymen}
\address
{School of Mathematics, Southampton University, Southampton SO17 1BJ,  England 
\emph{and} School of Mathematics, Alan Turing Building, Manchester University, Manchester M13 9PL, England}
\email{r.j.plymen@soton.ac.uk, roger.j.plymen@manchester.ac.uk}


\keywords{Endoscopy, characters, $p$-adic, $\SL_2$}
%\classno{20G05, 22E50}
%\date{\today}
\maketitle


\begin{abstract}  We consider the depth-zero supercuspidal $L$-packets of  $\SL_2(F)$ where $F$ is a non-archimedean local field of characteristic zero.   
We compare the explicit endoscopic character identities for $\SL_2(F)$ with  the classical character formulas of Sally-Shalika.   Our main result concerns the 
supercuspidal $L$-packet of size $4$.   For this $L$-packet, we show how the norm $1$ groups $H_1, H_2, H_3$ in the three quadratic extensions of $F$ play a crucial role in the 
endoscopic character identities for $\SL_2(F)$.
 

  
\end{abstract}

\tableofcontents

\section{Introduction}  We consider the depth-zero supercuspidal $L$-packets of  $\SL_2$ where $F$ is a non-archimedean local field of characteristic zero.  We compare the explicit endoscopic character identities for $\SL_2$ with  the classical character formulas of Sally-Shalika.

Our main result concerns the supercuspidal $L$-packet of size $4$.  
 For this $L$-packet, we show how the norm $1$ groups $H_1, H_2, H_3$ in the three quadratic extensions of $F$ play a crucial role in the 
endoscopic character identities for $\SL_2(F)$.
% With the aid of the classical character formulas of Sally-Shalika, we prove the endoscopic character identities.  

%\emph{Relation to the literature on $\SL_2$.} The reader may be tempted to think that ``everything about $\SL_2$ must already be known".  In order to counteract this viewpoint, we would like to orient our work in relation to the literature, and point out that the intersection of %our work with the literature is in fact rather minimal.    Lemma 7.1 in \cite{LL}.   7.1.2 in \cite{L}  .   Gelfand, Graev and Pyatetskii-Shapiro in \cite{GGS} pointed out that $\SL_2$ admits $L$-packets of cardinality $2$.   Shelstad established in \cite{Sh} that $\SL_2$ also admits %$L$-packets of cardinality $4$.   

The local field $F$ is a finite extension of $\Qp$.    We shall assume that $p > 2$. 
Let $q = q_F$ denote the cardinality of the residue field of $F$. Let $\val(x)$ denote the valuation of $x \in F^\times$.  Let $\fo_F$ denote the ring of integers of $F$,  let $\varpi$ be a uniformiser for $F$,  let $\fp_F=\varpi\fo_F$ be the maximal ideal of $\fo_F$, and let  $\varepsilon$ be a fixed nonsquare element in $\fo_F^\times$.  In that case, $\{1,\varpi,\varepsilon,\varepsilon\varpi\}$  is a set of representatives of $F^\times/(F^\times)^2$. 

Let $G = \SL_2(F)$, and let $\rZ(G)$ denote the center of $G$.  We shall focus on the elliptic torus
%For $\vartheta\in\{1,\varpi,\varepsilon,\varepsilon\varpi\}$, we define
\begin{equation} \label{eqn:tori}
T^\varepsilon:=
\left\{\left(
\begin{matrix}
a&b\cr
b\varepsilon&a
\end{matrix}\right)\;:\; a,b\in\fo_F, a^2-\varepsilon b^2=1
\right\}.
\end{equation}
This elliptic torus $T^\varepsilon$ is a representative of the single stable conjugacy class of unramified elliptic maximal $F$-tori in $\SL_2(F)$.   

We will define $f$ as follows:
\begin{equation} \label{eqn:definition-of-f}
f \colon T^\varepsilon \to \Z, \quad \quad 
\left(
\begin{matrix}
a&b\cr
b\varepsilon&a
\end{matrix}\right) \mapsto (-q)^{\val(b)} .
\end{equation}

We denote by $T^{\varepsilon}_1$ the pro-unipotent radical of $T^\varepsilon$, namely
\begin{equation} \label{T_1}
T^\varepsilon_1:=
\left\{\left(
\begin{matrix}
a&b\cr
b\varepsilon &a
\end{matrix}\right) \in T^\varepsilon \;:\; a \in 1 + \fp_F,\; b \in \fp_F \; \right\}.
\end{equation}

We will say, following \cite{ADSS}, that  $\gamma$ is \emph{near the identity} if $\gamma \in T_1^\varepsilon$,
and \emph{far from the identity} otherwise.  

%We recall the notions of Whittaker datum and generic representation for $\SL_2(F)$.    
We have the unipotent subgroup
\[
U(F) =\left\{ \left(
\begin{matrix}
1 & x \cr
0 & 1
\end{matrix}\right) : x \in F \right \} \subset \SL_2(F)
\]


 %A character $\theta : U(F) \to \C^\times$ is \emph{generic}  if it is nontrivial. 
 A Whittaker datum $\fw$  for $\SL_2(F)$  is a pair $(U, \theta)$ where   $\theta : U(F) \to \C^\times$ is a non-trivial character.     
 An irreducible smooth representation $(\pi, V)$ of $\SL_2(F)$  is called $\fw$-generic if the restriction of $\pi$ to $U(F)$ contains $\theta$.   
In this article,  we choose and fix a Whittaker datum $\fw$ for $\SL_2(F)$.   

If $G = \SL_2(F)$ then the Langlands dual group $\widehat{G}$ is $\PGL_2(\C)$ and the $L$-group ${}^LG$ is $\PGL_2(\C) \rtimes \W_F$.
In the context of this article, $\phi$ will denote a Langlands parameter
\[
\phi : \W_F \to {}^LG.
\]  
Let $S_\phi$  denote the centralizer in $\PGL_2(\C)$ of the image of $\phi$, and $\Irr(S_\phi)$ the set of isomorphism classes of irreducible representations of $S_\phi$. We have a bijective map
\begin{equation}
\iota_\phi  \colon \Pi_\phi(\SL_2(F))  \to \Irr(S_\phi)
\end{equation}
that sends the $\fw$-generic representation of $\SL_2(F)$ to the trivial character of $S_\phi$.   

Given  $s \in S_\phi$, we attach to $(\phi,s)$ the  virtual character    
\begin{equation} \label{eqn:virtual}                                                       
\Theta_{\phi, s} := \sum_{\pi\in\Pi_\phi(\SL_2(F)) } \trace(\iota_\phi(\pi))(s) \cdot \Theta_{\pi}.
\end{equation}
Let $\widehat{H}$ denote the identity component of the centralizer of $s$ in $\widehat{G}$. The \emph{extended endoscopic triple}, as defined in \cite[\S6.2]{T},  is 
        \[
        \fe(s): = (s, H, {}^L\eta).
        \]     Then $\phi$ will factorize as follows:
       \[
       \begin{CD}
       \W_F @> \phi^H >> {}^LH @> {}^L\eta >> {}^LG.
       \end{CD}
       \]
           
        
       The map $\phi^H$ will be called the  \emph{little parameter} associated to $\phi$.   This  leads to the definition of the \emph{stable character}.   
 The stable character $\sS\Theta_{\phi^H}$ is defined as 
 \[
 \sS\Theta_{\phi^H} := \Theta_{\phi^H, 1}.
 \] 
 Note that $H$ is an elliptic torus and the right-hand-side collapses to a single character of $H$. 
 
Then the  endoscopic character identity for $\SL_2$ is     
\begin{eqnarray}\label{endo}
 \sum_{\delta \in H(F)} \Delta[\fw,\fe(s)](\delta, \gamma) \sS\Theta_{\phi^H}(\delta) = \Theta_{\phi,s}(\gamma) 
\end{eqnarray}       
where $\Delta[\fw,\fe(s)]$ is the Langlands-Shelstad transfer factor.
 In the context of quasi-split connected reductive groups, our reference for the endoscopic character identity is equation (6.4) in the article by Taibi \cite{T}.       
  From a purely notational point of view, we prefer to switch $\gamma$ and $\delta$ in Taibi's equation (6.4).
 
 
    


 The explicit endoscopic character identities, for the depth-zero supercuspidal $L$-packets of $\SL_2$, are stated in Theorems \ref{A} and \ref{B}.     



% We call an $L$-parameter $\phi$ \textit{regular} (resp. \textit{non-singular}) if the representations in the $L$-packet  $\Pi_\phi(\SL_2(F))$ are regular (resp. non-singular) in the sense of \cite{Ka2,Ka3}. We note that the regular $L$-parameter in Theorem \ref{A} fits into the %framework of \cite{Ka2}, whereas the non-regular $L$-parameter in Theorem \ref{B} does not (but, being non-singular, fits into the more general framework of \cite{Ka3}).      

The characters of the discrete series for $\SL_2(F)$  were obtained by Sally-Shalika \cite{SS}.     In particular, for the depth-zero supercuspidal $L$-packets, the characters 
% For the unique depth-zero supercuspidal $L$-packet with non-regular parameter, the four characters $\Theta_1, \Theta_2, %\Theta_3, \Theta_4$ 
%of the unique depth zero supercuspidal $L$-packet with bob-regular parameter representations in the packet $\Pi_\phi$ 
appear in Table 3 of \cite{SS}.  They are given as locally integrable functions on the maximal torus of $\SL_2(F)$ and on the elliptic tori in $\SL_2(F)$.   Many years later, proofs were supplied by  Adler, DeBacker, Sally and Spice in \cite{ADSS}.   For the most part, we find it more convenient to refer to this paper  \cite{ADSS}. 

Let $E$ denote the unramified quadratic extension $F(\sqrt \varepsilon)$ and let $E^1$ denote the group of elements of norm $1$ in $E$.   We have a canonical isomorphism of elliptic tori:
\begin{eqnarray}
E^1 \to T^\varepsilon, \quad \quad a + b\sqrt \varepsilon \mapsto \left(
\begin{matrix}
a&b\cr
b\varepsilon&a
\end{matrix}\right)
\end{eqnarray}
and we shall often identify $E^1$ with $T^\varepsilon$.     A character $\psi$ of $E^1$ is \emph{quadratic} if $\psi^2 = 1$ and 
\emph{non-quadratic} if $\psi^2 \neq 1$.   The \emph{regular supercuspidal parameters} are defined by Kaletha \cite[Def. 5.2.3]{Ka2}.   
The group $\SL_2(F)$ admits a unique non-regular parameter, which features in Theorem \ref{B}.    


When we combine Theorems \ref{vch}  and  \ref{ENDO}, we obtain the following result.

%The endoscopic datum $\fe$ which appears in the following results is constructed in  \S7 and \S8.    

\begin{thm}\label{A}    Let $\gamma$ be a regular, semisimple element in $T^\varepsilon$.  Let $\Pi_\phi(\SL_2(F))$ be the depth-zero supercuspidal $L$-packet with regular $L$-parameter $\phi$.    
  %\begin{thm}  Let $\gamma$ be a regular semisimple element in $\SL_2(F)$.      We have 
 %Let $\gamma$ be a regular semisimple element in the elliptic torus $T^\varepsilon$ which is \textbf{far from the identity}.  
 Let $s$ in (\ref{endo}) denote the non-trivial element in $S_\phi$.  Then the endoscopic group $H$ is $E^1$.   The little parameter $\phi^H$ is the Langlands parameter of a non-quadratic depth-zero character
 $\psi$ of $H$.   
 %$ is the group of norm $1$ elements in the unramified quadratic extension $F(\sqrt \varepsilon)$.   
 We have  
  \begin{eqnarray*}\label{AA}
 \sum_{\delta \in H(F)}  \Delta[\fw, \fe(s)](\delta, \gamma) \sS\Theta_{\phi^H} (\delta) 
&=&- f(\gamma)(\psi(\gamma) + \psi(1/\gamma))
\end{eqnarray*}
and
\begin{eqnarray*}
\Theta_{\phi, s}(\gamma) &=&
\begin{cases} - \psi(\gamma) - \psi(1/\gamma) \quad  \text{when $\gamma$ is far from the identity} \\
-2f(\gamma) \quad \text{when $\gamma$ is near the identity} 
\end{cases}
  \end{eqnarray*}
%  where $\psi$ is a depth-zero character of $T^\varepsilon$ such that $\psi \neq \psi^{-1}$.  
%This equation is consistent with the  Sally-Shalika character formulas.
% we have   \begin{eqnarray*}  \Theta_{\phi, s}(\gamma) &=& -\psi(\gamma) - \psi(1/\gamma) \quad \textrm{far from the identity} \\  \Theta_{\phi, s}(\gamma) &=& -2f(\gamma) \quad \textrm{near the identity}    \end{eqnarray*}
%The endoscopic character identity is, in this case, consistent with the Sally-Shalika character formulas.
\end{thm}

These equations are consistent because $f(\gamma) = 1$ when $\gamma$ is far from the identity, and $\psi(\gamma) = 1$ when $\gamma$ is near the identity.   

The three quadratic extensions of $F$ are 
\[
E_1 = F(\sqrt \varepsilon),  \quad \quad E_2 = F(\sqrt \varpi), \quad \quad E_3 = F(\sqrt{\varepsilon \varpi}).
\]
Let $H_j$ denote the group of norm $1$ elements in $E_j$ with $j = 1,2,3$. 

The left-hand-side of the endoscopic character identity depends on the element $s$ in $S_\phi$.    In order to mark this dependence, we introduce the following definition:  
\[
\mathscr{E}(s): =  \sum_{\delta \in H}  \Delta[\fw, \fe(s)](\delta, -) \sS\Theta_{\phi^{H}} (\delta) 
\]
The value of $\mathscr{E}(s)$ at an element $\gamma$ will be denoted as follows: 
\[
\mathscr{E}(s:\gamma): =  \sum_{\delta \in H}  \Delta[\fw, \fe(s)](\delta, \gamma) \sS\Theta_{\phi^{H}} (\delta) 
\]


%The formula on the right-hand-side of (\ref{AA})  ``looks different" according as $%\gamma$ is far from the identity, or near the identity.   
The field $F$ admits a unique biquadratic extension field $K/F$ and the non-regular parameter $\phi$ factors through the Galois group $\Gal(K/F)$, see \S7.    Let $s_1, s_2$ and $s_3$ denote the images in $\PGL_2(\C)$ of the matrices
\[
\left(\begin{matrix}1&0\\0&-1\end{matrix}\right), \quad \left(\begin{matrix}0&1\\-1&0\end{matrix}\right),  \quad   \left(\begin{matrix}0&1\\1&0\end{matrix}\right).
\] 
We attach to $s_i$ the extended endoscopic 
triple $\fe(s_i) = (s_i, H_i, {}^L\eta_i)$ where ${}^L\eta_i : {}^LH_i \to {}^LG$ extends the embedding $\widehat{H_i} \to \widehat{G}$.      


  When we combine Theorems \ref{FAR-NON-REG}, 
\ref{NEAR-NON-REG}  and  \ref{abc}, we obtain the following result.   An interesting feature of this result is that all three endoscopic groups $H_1, H_2, H_3$ are involved.   

\begin{thm} \label{B}   Let $\gamma$ be a regular, semisimple element in $T^\varepsilon$.  Let $\Pi_\phi$ be the unique depth-zero supercuspidal $L$-packet with non-regular parameter $\phi$.  
Let $\psi_0$ denote the unique quadratic character of $T^\varepsilon$.   
%If $s = s_1$  then the endoscopic group $H$ is the same as in Theorem \ref{A}.   
%We have\begin{eqnarray*}s = s_1  & \implies & H = H_1,\quad  \fe = \fe_1\\s=s_2  & \implies & H = H_2, \quad \fe = \fe_2\\s = s_3  & \implies & H = H_3, \quad \fe = \fe_3\end{eqnarray*}
%The little parameter $\phi^H$ is the Langlands parameter of the unique quadratic character $\psi_0$ of $H$.   
 %$H$ is the group of norm $1$ elements in the unramified quadratic extension $F(\sqrt \varepsilon)$.   
We have 
\begin{eqnarray*}
\mathscr{E}(s_1 : \gamma) &=&  -2 f(\gamma) \psi_0(\gamma) \\
\mathscr{E}(s_2 : \gamma) &=& 0\\
\mathscr{E}(s_3:\gamma) & = & 0
\end{eqnarray*}
and 
\begin{eqnarray*}
\Theta_{\phi, s_1} &=&
\begin{cases} 
-2\psi_0(\gamma) \quad  \text{when $\gamma$ is far from the identity}\\
 -2f(\gamma)   \quad \text{when $\gamma$ is near the identity} 
\end{cases}\\
\Theta_{\phi, s_2}  &=& 0 \\
\Theta_{\phi, s_3} &=& 0.
%\Theta_{\phi, s_1} &=& 0 \quad \text{when $\gamma$ is near the identity}\\
%\Theta_{\phi, s_3} &=& -2f(\gamma)  \quad \text{when $\gamma$ is near the identity}
\end{eqnarray*}
%where $\psi_0$ is the unique quadratic character of $T^\varepsilon$.    
\end{thm} 



%We will use `Theorem 15.2' to refer to Theorem 15.2 in \cite{ADSS}.   Unfortunately, when $\gamma$ is near the identity,
%\[Theorem \; 15.2 \implies \begin{cases} \Theta_{\phi, s_1} = 0\\\Theta_{\phi, s_2} = 0\\\Theta_{\phi, s_3} = -2f(\gamma)\end{cases}  \]
%For a proof of this implication, see \S5.1.     The conclusion is inescapable: Theorem 15.2 is false.   We note that Theorem 15.2 occurs already in the 1968 article  of Sally-Shalika, see the 7th (incorrect) formula in Table 3 of \cite{SS}.
 

%From the Sally-Shalika character formulas, near the identity,  we have \footnote{there is %an unexplained discrepancy here}\begin{eqnarray*} \Theta_{\phi, s_3}(\gamma) &=& %-2f(\gamma) \end{eqnarray*}\end{thm}and $\Theta_{\phi, s_1}(\gamma) = \Theta_{\phi, %s_2}(\gamma) = 0$.
%Consistency in the formulas in Theorem \ref{B}  is due to Lemma \ref{farfrom}.   
%In \cite{FKS}, the authors Fintzen-Kaletha-Spice, in a much more general context,  consider a supercuspidal parameter $\phi : \bW_F \to {}^LG$ and let $S_\phi = \Cent(\phi, \widehat{G})$.
%There is an associated torus $S$ defined over $F$, a character $\theta : S(F) \to \C^\times$ and a canonical exact sequence
%\[1 \to \widehat{S}^\Gamma \to S_\phi \to \Omega(S,G)(F)_\theta \to 1\]where $\Omega(S,G)$ is a subgroup of the automorphism group of $S$ that is identified with
%the Weyl group coming from $G$ in a natural way, and $\Gamma$  is the absolute Galois group of $F$.     When $\phi$ is regular, $\Omega(S,G)_\theta$ is trivial; hence $S_\phi \simeq \widehat{S}^\Gamma$.  The authors prove
%the endoscopic character identities for all $s \in \widehat{S}^\Gamma \subset S_\phi$ , subject to certain conditions on the base field $F$ .
%Let $E^1$ denote the group of norm $1$ elements in the unramified quadratic extension $F(\sqrt \varepsilon)$.   In the context of our Theorem \ref{A}, we have 
%\[S = E^1, \quad \theta = \psi, \quad \Omega_\theta = \Z/ 2\Z, \quad  \widehat{S}^\Gamma = \{1, s_1\}
%\]where $s_1$ is the image in $\PGL_2(\C)$ of the matrix $\left(\begin{smallmatrix}1&0\\0&-1\end{smallmatrix}\right)$.   This would enable the authors to re-create Theorem \ref{A} and Theorem \ref{B} only in the case $s = s_1$.   
% Now, by (\ref{STheta}) and Lemma \ref{Delta3} we have \[\sum_{\gamma^H} \Delta[\fw, \fe](\gamma^H, \gamma) S\Theta_{\phi^H}(\gamma^H) = -2f(\gamma) \psi_0(\gamma^H)
%-2f(\gamma) & \text{when $\gamma$ is near the identity}\\
%-2\psi_0(\gamma^H) & \text{when $\gamma$ is far from the identity}\]and by Theorem \ref{FAR} and Theorem \ref{NEAR-NON-REG} we have\[\Theta_{\phi, s_1} =\begin{cases}-2\psi_0(\gamma^H) & \text{when $\gamma$ is far from the identity}\\
%-2f(\gamma)   & \text{when $\gamma$ is near the identity}\\\end{cases}\]so we have consistency.   However, in \S5.1 we show that\begin{equation}\text{Theorem 15.2 implies} \quad \Theta_{\phi, s_1} = 0  \quad \text{when $\gamma$ is near the identity}
%\end{equation}and it is here that we have inconsistency.  
%\emph{On Theorem 15.2 in \cite{ADSS}}.   This theorem is incompatible with Proposition 13.14 in \cite{ADSS} and with the endoscopic character identity, see \S5.    
   %We have used Proposition 13.14 in \cite{ADSS} instead.   

\emph{On the issue of uniqueness}.   In the context of a regular parameter $\phi$, 
the uniqueness of the map $\iota_\phi \colon \Pi_\phi \to \Irr(S_\phi)$ is clear. For the $\fw$-generic element in $\Pi_\phi$ is sent to the trivial character of $S_\phi$, whereas the other element in $\Pi_\phi$ is sent to the non-trivial character of $S_\phi$.  

In the context of the non-regular parameter $\phi$,  the left-hand-side of (\ref{endo}) depends, in its definition, on the map $\iota_\phi\colon \Pi_\phi \to \Irr(S_\phi)$.  The  endoscopic character identity selects a unique map
 $\iota_\phi : \Pi_\phi \to \Irr(S_\phi)$, see \S4.   

%\emph{Relation to the recent article by Fintzen-Kaletha-Spice \cite{FKS}}. 
%We would like to make the following observation.   
%Fintzen-Kaletha-Spice work in a much more general context.
% They work in the vastly more general context of supercuspidal representations $\pi_{(S, \theta)}$ indexed by tame 
%elliptic pairs $(S, \theta)$ where $\theta : S(F) \to \C^\times$ is a character.  
%In theorem 4.4.2 in \cite{FKS} the authors give a highly elaborat fprmula for the vitual caharcetr $
%According to \cite[p.2244]{FKS} there is a canonical exact sequence\[1 \to \widehat{S}^\Gamma \to S_\phi \to \Omega(S,G)_\theta \to 1\]
%where $\Gamma$ is the absolute Galois group $\Gal(\overline{F}/F)$,   $\Omega(S,G)$  is  the Weyl group of $G$, and $\Omega(S,G)_\theta$ is the  stabilizer of $\theta$.  When $\phi$ is regular, so that $\Omega(S,G)_\theta$ is  trivial, this sequence becomes an %isomorphism \[\widehat{S}^\Gamma \simeq S_\phi.\]This is the situation in Theorem \ref{A}.     
%In \cite[Theorem 4.4.2]{FKS}, the authors give an  elaborate formula for the virtual character $\Theta_{\phi, s}$ provided that 
%$s \in \widehat{S}^\Gamma$.  In our context with $S = E^1$ and $\theta = \psi$ or $\psi_0$, this means that $s = s_1$.      
%It follows that our formulas for $\Theta_{\phi, s_1}$ are a special case of their formula. However, our formulas for 
%$\Theta_{\phi, s_2}$ and $\Theta_{\phi, s_3}$ in Theorem \ref{B} are not a special case, since $s_2$ and $s_3$ lie outside the fixed set $\widehat{S}^\Gamma$. 
  
We discovered a crucial typo in the statement of Theorem 15.2 in \cite{ADSS}; the corrected statement appears in \S5.  

  \textbf{Acknowledgement}. We would like to thank Loren Spice for a valuable exchange of emails, Shaun Stevens for valuable discussions  and Tasho Kaletha for 
 sending us constructive criticism. 
 



 
   
   
\section{Elliptic tori}   
\subsubsection{Elliptic tori} There are two tori relevant for us:
 Let $\varepsilon$ be a fixed nonsquare element in  $\fo_F^{\times}$, then

\begin{equation} \label{eqn:Torus-Tepsilon}
T^\varepsilon:=
\left\{\left(
\begin{matrix}
a&b\cr
b\varepsilon &a
\end{matrix}\right)\;:\; a,b\in\fo_F, a^2-\varepsilon b^2=1
\right\}
\end{equation}
and 
\begin{equation} \label{eqn:Torus-Tepsilon-varpi}
T^{\varepsilon, \varpi} :=
\left\{\left(
\begin{matrix}
a&b\varpi^{-1} \cr
b\varepsilon \varpi &a
\end{matrix}\right)\;:\; a,b\in\fo_F, a^2-\varepsilon b^2=1
\right\}.
\end{equation}
These two tori admit the same splitting field, namely $F(\sqrt \varepsilon)$ and are therefore stably conjugate, or, equivalently,
$\GL_2(F)$-conjugate,    thanks to the following identity in $\GL_2(F)$:
\begin{equation}\label{ggg}
\left(\begin{array}{cc}
1 & 0\\
0 & \varpi
\end{array}
\right)
\left(\begin{array}{cc}
a & b\\
\varepsilon b & a
\end{array}
\right)
\left(\begin{array}{cc}
1 & 0\\
0 & \varpi^{-1}
\end{array}
\right)
 = 
\left(\begin{array}{cc}
a & \varpi^{-1}b\\
\varepsilon \varpi b & a
\end{array}\right).
\end{equation}
There is a  single stable conjugacy class of unramified elliptic maximal $F$-tori,  represented by $T^\varepsilon$.   It splits into two 
distinct $G$-conjugacy classes, represented by $T^\varepsilon$ and $T^{\varepsilon, \varpi}$.   

Let $E^1$ denote the group of norm $1$ elements in $E$.   We have, following \cite[\S3.1]{ADSS}, canonical isomorphisms of $F$-tori:
\[
T^\varepsilon \to E^1, \quad \left(
\begin{matrix}
a&b\cr
b\varepsilon &a
\end{matrix}\right)
\mapsto a + b\sqrt\varepsilon
\]
\[
T^{\varepsilon, \varpi} \to E^1, \quad \left(
\begin{matrix}
a&b\varpi^{-1}\cr
b\varepsilon \varpi &a
\end{matrix}\right)
\mapsto a + b\sqrt\varepsilon.
\]
Given $a + b\sqrt \varepsilon \in E^1$ we shall, following \cite[\S3.1]{ADSS} write 
\[
\mathrm{Im}_\varepsilon(a + b\sqrt \varepsilon) = b.
\]  
We have
\begin{eqnarray}\label{gdot}
\mathrm{Im}_\varepsilon (\g \cdot \gamma) =  \mathrm{Im}_\varepsilon (\gamma),
\end{eqnarray}
for every $\gamma\in T^{\varepsilon}$, where 
\begin{equation} \label{eqn:g}
\mathbf{g}: = \left(
\begin{array}{cc}
1 & 0\\
0 & \varpi
\end{array}
\right) \in \GL_2(F)
\end{equation}
and $\mathbf{g} \cdot \gamma: = \mathbf{g} \gamma \mathbf{g}^{-1}$.
The equation \eqref{ggg} can be rephrased as the following statement:
\begin{equation} \label{eqn:tori}
\g \cdot T^\varepsilon =  \g T^\varepsilon \g^{-1} = T^{\varepsilon, \varpi}.
\end{equation}
We also note that
 \[
 \sgn_\varepsilon(x) = (-1)^{\val(x)}
 \]
 for all $x \in F$, see (2.1) in \cite {ADSS}. Then we have, using (\ref{gdot}),
 \begin{eqnarray}\label{SGN}
 \sgn_\varepsilon(\varpi^{-1} \mathrm{Im}_\varepsilon(\g \cdot \gamma))  &=& \sgn_\varepsilon(\varpi^{-1} \mathrm{Im}_\varepsilon(\gamma)) \\
  &=& \sgn_\varepsilon(\varpi^{-1}) \sgn_\varepsilon( \mathrm{Im}_\varepsilon(\gamma)) \nonumber \\
 &=& (-1)^{\val(\varpi^{-1})}\sgn_\varepsilon  (\mathrm{Im}_\varepsilon(\gamma)) \nonumber\\
 &=&- \sgn_\varepsilon( \mathrm{Im}_\varepsilon(\gamma)). \nonumber
 \end{eqnarray}


\subsubsection{Weyl discriminant} In \cite[Definition 3.7]{ADSS}, the Weyl discriminant $D_G$ is defined as follows:  
\[
D_G \left(
\begin{array}{cc}
a & b\\
c & d
\end{array}
\right)
= (a + d)^2 - 4.
\]
%Let $T$ be the diagonal subgroup of $\SL_2(F)$.   The usual definition of  the Weyl discriminant $D^G$  is the following. If $\gamma \in T$ then \begin{equation} \label{eqn:discriminant}
%D^G(\gamma) = | \Pi_{\alpha} (\alpha(\gamma) - 1)|^{1/2} \end{equation} where the product is over all roots of $T$ in $G$.  

%If\[\gamma = \left( \begin{array}{cc}\lambda & 0\\0 & \lambda^{-1}\end{array}\right)\] then\[D_G(\gamma) = (\lambda + \lambda^{-1})^2 -4  = (\lambda - \lambda^{-1})^2.\]
%We also have\begin{eqnarray*}D^G(\gamma) &=&| \Pi_{\alpha} (\alpha(\gamma) - 1)|^{1/2}\\&=& |(\lambda^2 - 1)(\lambda^{-2} - 1)|^{1/2}\\&=& |\lambda(\lambda - \lambda^{-1})(\lambda^{-1} - \lambda)\lambda^{-1}|^{1/2}\\
%&=& |\lambda - \lambda^{-1}|.\end{eqnarray*}It follows that\begin{eqnarray}\label{Weyl}|D_G(\gamma)|^{1/2} = D^G(\gamma).\end{eqnarray}The term $|D_G(\gamma)|^{1/2}$ occurs repeatedly in the denominator of the character formulas in \cite{ADSS}.   Thanks to equation (\ref{Weyl}), we shall write this denominator as $D^G(\gamma)$. 
  
Now
\[
\gamma = \left(
\begin{matrix}
a&b\cr
b\varepsilon &a
\end{matrix}\right) \in T^\varepsilon
\implies D_G(\gamma) =  4a^2 - 4 = 4 \varepsilon b^2
\]
so that
\[
|D_G(\gamma)|^{1/2}  = |b|.
\]

\subsubsection{The function $f$} We set
\begin{equation} \label{eqn:fgamma}
f(\gamma):= \frac{\sgn_\varepsilon(\mathrm{Im}_\varepsilon(\gamma))}{|D_G(\gamma)|^{1/2}} .
\end{equation}
Note that 
\[
\gamma \quad  \textrm {is regular} \iff b \neq 0.
\]
In that case, we have
\begin{eqnarray*} 
f(\gamma)&=& \frac{\sgn_\varepsilon(b)}{|b|}\\
&=& (-1)^{\val(b)} q^{\val(b)}\\
&=& (-q)^{\val(b)}.
\end{eqnarray*}

It is worth noting that $f$ is invariant under the Galois group $\Gal(F(\sqrt \varepsilon)/F)$.  If $c = a + b\sqrt \varepsilon$ then we have
\begin{equation}
f(\overline{c}) = f(a - b\sqrt \varepsilon) = (-q)^{\val(-b)} = (-q)^{\val(b)} = f(c).
\end{equation}\label{fGalois}
%Note that $D_G(\g \cdot \gamma) = D_G(\gamma)$ and so we have \[f(\g\cdot \gamma) = f(\gamma).\]

The following Lemma is proved in \cite[p.65]{ADSS}.    We offer a short proof.

\begin{lem}\label{farfrom}  Let $\gamma$ be a regular semisimple element 
in the elliptic torus $T^\varepsilon$ which is far from the identity.   Then we have $f(\gamma) = 1$.
\end{lem}
\begin{proof} We recall the pro-unipotent radical $T^\varepsilon_1$ of $T^\varepsilon$, see (\ref{T_1}).  If $\gamma$ is far from the identity, then we have 
$\gamma \in T^\varepsilon \setminus T^\varepsilon_1$ and so $b \in \fo_F \setminus \fp_F$.  Therefore $\val(b) = 0$. 
\end{proof}


\subsubsection{Elliptic torus in $\SL_2(\kF)$}   The quotient $T^\varepsilon/T_1^\varepsilon$ is the unique maximal elliptic $\kF$-torus $\bbT^\epsilon(\kF)$ in $\SL_2(\kF)$:
\begin{equation}
\bbT^\epsilon (\kF)=\left\{ \left(\begin{matrix}
   a & b\\
 \epsilon b & a
 \end{matrix}
 \right)\,:\, \text{$a,b\in \kF$ and $a^2-\epsilon b^2=1$}\right\}
 \end{equation} 
 where $\kF$ is the residue field of $F$.   

  %We now proceed to construct the endoscopic datum $\fe$.   First, we need the quadratic character of the unramified elliptic torus.   
  
The group $E^1$ of norm one elements in $E = F(\sqrt \varepsilon)$ admits a unique quadratic character $\psi_0$ of depth zero which therefore factors through $\bbT^\epsilon$.  
An explicit formula for $\psi_0$ (which we shall not need) is given in \cite[\S 9.2]{ADSS}.   

   

\section{Depth-zero supercuspidal representations of $\SL_2(F)$}  
Let $I$ denote a line-segment (chamber) in the reduced building  (tree) of $\SL_2(F)$.   The ends of the interval $I$ will be denoted $x_0$, $x_1$ .   The stabilizer of $x_i$ will be denoted $G_{x_i}$, for $i=0,1$. We have  $G_{x_0}=\SL_2(\fo_F)$, and  $G_{x_1}$ is the conjugate of $G_{x_0}$ under the element $\mathbf{g}$ of $\GL_2(F)$ defined in \eqref{eqn:g}.
The reductive quotients of $G_{x_0}$ and $G_{x_1}$ are both  isomorphic to  $\SL_2(\kF)$,  and $\{G_{x_0},G_{x_1}\}$ is a set of representatives of the maximal parahoric subgroups of $\SL_2(F)$.


Let $\psi$ be a character of $\bbT^\epsilon(\kF)$ such that $\psi \ne \psi^{-1}$.  The Deligne-Lusztig virtual character $R_{\bbT^\epsilon,\psi}^{\SL_2}$ is irreducible,  and  we denote by $|R_{\bbT^\epsilon,\psi}|$  its inflation
to $G_{x_0}$ and $G_{x_1}$. We put
\begin{equation} \label{eqn:reg}
\pi^+=\pi(\bbT^\epsilon,\psi):= c-\Ind_{G_{x_0}}^G |R_{\bbT^\epsilon,\psi}|\quad\text{and}\quad\pi^-:= c-\Ind_{G_{x_1}}^G |R_{\bbT^\epsilon,\psi}|.
\end{equation}
The representations $\pi^+$ and $\pi^-$  are depth-zero regular supercuspidal  irreducible representations of $G$ in the terminology of \cite{Ka2}.

Let $\psi_0$ be the quadratic character of $\bbT^\epsilon(\kF)$. The Deligne-Lusztig virtual character $R_{\bbT^\epsilon,\psi_0}^{\SL_2}$  decomposes as 
\begin{equation}
R_{\bbT^\epsilon,\psi_0}^{\SL_2}=-R_{\bbT^\epsilon,\psi_0}^+ - R_{\bbT^\epsilon,\psi_0}^-,
\end{equation}
where $R_{\bbT^\epsilon,\psi_0}^\pm$ are irreducible cuspidal representations of $\SL_2(\kF)$ of dimension $\frac{q-1}{2}$ (see \cite[Table~12.1]{DM}). We choose the signs as in \cite[(9.3)]{ADSS}.

We denote by $|R_{\bbT^\epsilon,\psi_0}^\pm|$  the inflations of $R_{\bbT^\epsilon,\psi_0}^\pm$ to $G_{x_0}$ and $G_{x_1}$.
The four compactly induced representations $\cInd_{G_{x_i}}^G |R_{\bbT^\epsilon,\psi_0}^\pm|$, $i=0,1$, are depth-zero supercuspidal  irreducible representations of $G$, they are called ``exceptional" in \cite{ADSS}, and are non-singular non-regular  in the terminology of  \cite{Ka3} (see also \cite{Au2}). We write
\begin{equation} \label{eqn:nonsingular}
\pi^+:=\cInd_{G_{x_0}}^G |R_{\bbT^\epsilon,\psi_0}^+|\quad\text{and}\quad
\pi^-:=\cInd_{G_{x_1}}^G |R_{\bbT^\epsilon,\psi_0}^-|.
\end{equation}
We define
 \begin{equation*}
 (\g \cdot \pi^+)(x):= \pi^+(\g^{-1}x\g)\quad\text{and}\quad
 \g \cdot x:= \g x\g^{-1},
 \end{equation*}
 so that 
 \[
 (\g \cdot \pi^+)(x) = \pi^+(\g^{-1} \cdot x).
 \]
 We set 
 \begin{equation} \label{eqn:packet}
\pi_1 = \pi^+, \quad \pi_2 = \pi^-,\quad \pi_3 = \g \cdot \pi^+, \quad \pi_4 = \g\cdot \pi^- .
\end{equation}

\section{Virtual characters far from the identity: the non-regular parameter}  We recall that  $G = \SL_2(F), \, \widehat{G} = \PGL_2(\C) = \PSL_2(\C)$.   
Let $E/F$ be the (unramified) quadratic extension of $F$ containing $\sqrt \varepsilon$ and let $K/F$ be the biquadratic extension of $F$ containing $\sqrt \varepsilon$ and $\sqrt \varpi$.    
 %   Let $E/F$, $E'/F$ be distinct quadratic extensions, with $E/F$ unramified.   Let $\sigma$ (resp. $\tau$) be the non-trivial element in $\Gal(E/F)$ (resp. $\Gal(E'/F)$).     
% Let $K/F$ denote the biquadratic extension comprising $E/F$ and $E'/F$.   Let $s, \n$ be a pair of commuting semisimple elements in $\widehat{G}$ of order $2$.   
For the Galois group  we have $\Gal(K/F) = \{1,\sigma, \tau, \sigma \tau\}$ with $\sigma \sqrt \varepsilon = - \sqrt \varepsilon$ and $\tau \sqrt \varpi = - \sqrt \varpi$.   

  The Langlands parameter 
 \begin{eqnarray}\label{biquad}
 \phi \colon \bW_F \to \PGL_2(\C) 
 \end{eqnarray}
  is now defined   by the condition that $\phi$ factors through $\Gal(K/F)$ and satisfies the equations 
 \begin{eqnarray}\label{pd1}
  \phi(\sigma) &=& s_2, \quad   \phi(\tau) = s_1.
  \end{eqnarray}  
  
We consider the associated supercuspidal $L$-packet  
\[
\Pi_\phi(\SL_2(F)) = \{\pi_1, \pi_2, \pi_3. \pi_4\},
\]
and let $\Theta_j := \Theta_{\pi_j}$ with $1 \leq j \leq 4$. 

%The characters of the representations $\pi^+$ and $\pi^-$ (attached to $G_{x_0}$) were computed  in \cite{SS}. 

\begin{thm}\label{FAR}  Let $\gamma$ be a regular semisimple element 
in the elliptic torus $T^\varepsilon$ which is far from the identity. Then we have the character formulas
 %up to a harmless factor \footnote{cf. Kaletha, comment after formula for $\Delta(c,c)$  \cite[p.27]{Kal1}} of modulus $2$, we can reconcile the %character formula \ref{ADSSch} with the formula coming from endoscopic transfer.   The two formulas are, respectively, 
 \begin{eqnarray*}
\Theta_1(\gamma)  &=&  - \psi_0(\gamma)\\
 \Theta_2(\gamma)  &=& -\psi_0(\gamma)\\
  \Theta_3(\gamma)  &=&  0\\
 \Theta_4(\gamma)  &=& 0.
 \end{eqnarray*}    
 \end{thm} 
 \begin{proof}   We refer to equations $(*)$ and $(**)$ on p.65 of \cite{ADSS}.  It follows immediately that 
  \[
 \Theta_{\pi^+}(\gamma) =  \Theta_{\pi^-}(\gamma)  = - \psi_0(\gamma).
 \]
 We have, for all $\gamma \in T^\varepsilon$, 
 \begin{eqnarray*}
 \Theta_{\pi_3}(\gamma) &=& \Theta_{\g \cdot \pi^+}(\gamma)\\
 &=& \Theta_{\pi^+}(\g^{-1}\cdot \gamma)\\
 &=& 0
 \end{eqnarray*}
 since $\Theta_{\pi^+}$ vanishes on $\g^{-1} \cdot T^\varepsilon$ by \cite[Theorem 15.1]{ADSS}.   
 
Similarly, we have, for all $\gamma \in T^\varepsilon$, 
\begin{eqnarray*}
 \Theta_{\pi_4}(\gamma) &=& \Theta_{\g \cdot \pi^-}(\gamma)\\
 &=& \Theta_{\pi^-}(\g^{-1}\cdot \gamma)\\
 &=& 0
 \end{eqnarray*}
 since $\Theta_{\pi^-}$ vanishes on $\g^{-1} \cdot T^{\varepsilon}$ by \cite[Theorem 15.1]{ADSS}.
 \end{proof}
 
 Consider the $L$-packet of $\SL_2(F)$ given by
\begin{equation}
\Pi_{\phi}(\SL_2(F)) = \{\pi_1, \pi_2, \pi_3, \pi_4\}
%=\left\{\cInd_{K}^G|R_{\bbT_\varepsilon,\psi}^\pm|,\cInd_{K'}^G|R_{\bbT_\varepsilon,\psi}^\pm|\right\}.
\end{equation}
where the representations $\pi_i$ are depth-zero supercuspidal irreducible representations of $G$ (see \eqref{eqn:packet}).
We  suppose that $\pi_1$ is $\fw$-generic.   

Let  $\Theta_1$,  $\Theta_2$, $\Theta_3$, $\Theta_4$ be the distribution characters of $\pi_1$, $\pi_2$, $\pi_3$, $\pi_4$. 
%where\[S_{\varphi} = Q/\{1,-1\}= \{1,a,b,c\}\]the non-cyclic group of order $4$. 
%and the matrix of coefficients (the character table) is (as always) invertible.
%\section{Endoscopic triples for $\SL_2$}    

Given $A \in \GL_2(\C)$ let $A_*$ denote the image of $A$ in $\PGL_2(\C)$.   Let
       \begin{equation} \label{eqn:si}
       s_1 = \left(
       \begin{array}{cc}
       1 & 0\\
       0 & -1
       \end{array}
       \right)_*
             \quad s_2 = \left(\begin{array}{cc}
       0 & 1\\-1 & 0
       \end{array}
       \right)_*
       \quad 
       s_3 = \left(
       \begin{array}{cc}
       0 & 1\\
       1 & 0
       \end{array}
       \right)_*.
       \end{equation}  
 We have
       \[
       S_\phi = \{1, s_1, s_2, s_1s_2 \}.
       \]
We will enumerate the characters of $S_\phi$ in the following way:
       \[
       \begin{tabular}{c|cccc}
       & $1$ & $s_1$ & $s_2$ & $s_1s_2$\\
       \hline
       $\rho_1$ & $1$ & $1$ & $1$ & $1$\\
       $\rho_2$ & $1$ & $1$ & $-1$ & $-1$\\
       $\rho_3$ & $1$ & $-1$ & $1$ & $-1$\\
       $\rho_4$ & $1$ & $-1$ & $-1$ & $1$
       \end{tabular}
       \]

Given 
\begin{equation} 
\iota \colon\Pi_\phi \to \Irr(S_\phi)
\end{equation}
 we attach to $(\phi,s)$ the  virtual character    
\begin{equation} \label{eqn:virtual}                                                       
\Theta_{\phi, s} := \sum_{\pi\in\Pi_\phi(\SL_2(F)) } \iota(\pi)(s) \cdot \Theta_{\pi}.
\end{equation}
   
The first constraint is that $\iota$ must assign the trivial character 
$\rho_1$ to the $\fw$-generic representation $\pi_1$.   The second important constraint is the endoscopic character identity.  
 These constraints determine the map $\iota$ uniquely.  The defining equation is
\begin{equation} 
\iota(\pi_j) : = \rho_j \quad \text{ with $j = 1,2,3,4$.}
\end{equation}
In that case, we have
\begin{equation} \label{eqn:endos}
\Theta_{\phi, s} =\begin{cases} \Theta_1 + \Theta_2 + \Theta_3 + \Theta_4 &\text{if $s = 1$}\\
 \Theta_1 + \Theta_2 - \Theta_3 - \Theta_4   &\text{if $ s = s_1$}\\
\Theta_1 - \Theta_2 + \Theta_3 - \Theta_4.   &\text{if $s = s_2$}\\
\Theta_1 - \Theta_2 - \Theta_3 + \Theta_4    &\text{if $s = s_3$}
\end{cases}.
\end{equation}


%Similarly, $\Theta_4 = 0$ on the elliptic torus $T^\varepsilon$, since $\pi_4$ is conjugate via the same element $g$ to $\pi^-$.   
 \begin{thm}\label{FAR-NON-REG}  Let $\gamma$ be a regular semisimple element 
in the elliptic torus $T^\varepsilon$ which is far from the identity. Then we have
  \begin{eqnarray*}
\Theta_{\phi,1\hphantom{s}}(\gamma)  &=&  -2 \psi_0(\gamma)\\
 \Theta_{\phi, s_1}(\gamma)  &=& -2\psi_0(\gamma)\\
  \Theta_{\phi, s_2}(\gamma)  &=&  0\\
 \Theta_{\phi, s_3}(\gamma)  &=& 0.
 \end{eqnarray*}    
 \end{thm} 
 \begin{proof}   This follows readily from Theorem \ref{FAR}.   
 \end{proof}
 
 
 
   \section{Virtual characters near the identity: the non-regular parameter}  
 
 The sign character $\sgn_\varepsilon \colon F^\times \to \{\pm 1\}$ is determined by the short exact sequence
  \[
  1 \to \rN_{E/F}E^\times \to F^\times \to \Gal(E/F) \to 1.
  \]
  The notation $\mathrm{Im}_\varepsilon$ has already been defined in \S3.1.        
  We recall the definition of $f(\gamma)$ from \eqref{eqn:fgamma}.
    
    
%We will use the Cayley transform defined on an open subset of $\fg$ by\[c\colon X\mapsto \frac{1+X/2}{1-X/2}.\]We recall from \cite[Lemma~5.5(h)]{ADSS} that $|D_\fg(Y)|=|D_G(c(Y))|$ when $Y\in\fg_{0+}$.
 \begin{thm}\label{NEAR-NON-REG}  Let $\gamma$ be a regular semisimple element 
in the elliptic torus $T^\varepsilon$ which is near the identity.  Then we  have 
 \begin{eqnarray*}
 \Theta_{\phi, 1}(\gamma) &=& - 2\\
 \Theta_{\phi, s_1}(\gamma) &=& -2f(\gamma)\\
 \Theta_{\phi, s_2}(\gamma) &=& 0\\
 \Theta_{\phi, s_3}(\gamma) & = & 0
 \end{eqnarray*}
 \end{thm}
 \begin{proof}  According to Theorem 15.2 in \cite{ADSS}, we have \footnote{we have made a surely necessary correction to \cite[Theorem 15.2]{ADSS}, deleting the $\pm$ inside the curly bracket} 
 \[
 \Theta_{\pi^\pm}(\gamma) = \frac{1}{2} \left\{H(\Lambda', F_{\theta'}) \frac{\sgn_{\theta'}(\eta^{-1} \mathrm{Im}_{\theta'}(\gamma))}{|D_G(\gamma)|^{1/2}} - 1\right\}
 \]  
 
With $\gamma \in T^\varepsilon, \theta' = \varepsilon, \eta = 1$ we have
 \begin{eqnarray*}
 \Theta_1(\gamma) & = & \Theta_2(\gamma) \\
 & = & ( - f(\gamma) - 1)/2
 \end{eqnarray*}
 
 With $\gamma \in T^\varepsilon, \g \cdot \gamma \in T^{\varepsilon, \varpi}, \theta' = \varepsilon, \eta = \varpi$ we have
 \begin{eqnarray*}
 \Theta_4(\gamma) & = & \Theta_3(\gamma) \\
 &=&  \Theta_1(\g \cdot \gamma)\\ 
 &=& (f(\gamma) - 1)/2
 \end{eqnarray*}
 and the result follows immediately.
 \end{proof}



 
\section{Virtual characters: regular parameters}  We recall from \eqref{eqn:reg} the definition of  the depth-zero supercuspidal representation
\[
\pi^+:= \pi(T^\varepsilon, \psi).
\]
The $L$-packet is then 
\[
\Pi_\phi = \{\pi^+, \pi^-\}
\]
with $\pi^- = \g \cdot \pi^+$.




\begin{thm}\label{vch}  Let $s$ be the non-trivial element in $S_\phi$.   Let $\gamma$ be a regular semisimple element 
in the elliptic torus $T^\varepsilon$. Then we have
\begin{eqnarray*}
\Theta_{\phi, s} (\gamma) & = &
\begin{cases}
 -\psi(\gamma) - \psi(1/\gamma) \quad  \text{when $\gamma$ is far from the identity} \\
-2f(\gamma) \quad \text{when $\gamma$ is near the identity} 
\end{cases}
\end{eqnarray*}
\end{thm}
\begin{proof}   First case: $\gamma$ is far from the identity.    The equation 
\[
\Theta_{\pi^+}(\gamma) = -\psi(\gamma) - \psi(1/\gamma)
\]
 is a re-statement of Theorem 14.14 in \cite{ADSS}.   We have simplified this statement  by using Lemma \ref{farfrom} and \cite[Lemma 4.2]{ADSS}.   
 According to \cite[Theorem 14.14]{ADSS} we have
 \[
 \Theta_{\pi^-}(\gamma) = 0.
 \]
 
 We recall that $S_\phi = \{1,s\} = \Z /2\Z$.  Denote the trivial character of $S_\phi$ by $\rho_0$ and the non-trivial 
character by $\rho_1$.   Assume that $\pi$ is 
$\fw$-generic so that the parametrization of the $L$-packet is 
\[
\iota \colon \Pi_\phi \to \Irr(S_\phi), \quad \quad \pi^+ \mapsto \rho_0, \quad \pi^- \mapsto \rho_1.
\]
 
 In that case, we have
\begin{eqnarray*}                                                       
\Theta_{\phi, s}(\gamma) &=& \sum_{\pi\in\Pi_\phi(\SL_2(F)) } \trace(\iota(\pi))(s) \cdot \Theta_{\pi}\\
&=& \rho_0(s) \Theta_\pi(\gamma) + \rho_1(s)\Theta_{\pi^-}(\gamma)\\
&=& \Theta_{\pi^+}(\gamma)\\
&=& -\psi(\gamma) - \psi(1/\gamma).
\end{eqnarray*}

Second case: $\gamma$ is near the identity.    In the terminology of \cite[Definition 14.1]{ADSS}, we are in the \emph{unramified case}.  According to \cite[Theorem 14.20]{ADSS} we have
 \[
 \Theta_{\pi^+}(\gamma) = c_0(\pi) - f(\gamma)
 \]
 where, by \cite[Definition 14.17]{ADSS}, the constant term $c_0(\pi)$ is 
 \[c_0(\pi) = -1
 \]
 and so
\[
\Theta_{\pi^+}( \gamma) = -1 - f(\gamma).
\]

Note that $\gamma \in T^\varepsilon \implies \g \cdot \gamma \in T^{\varepsilon, \varpi}$.   
With $\gamma \in T_1^{\varepsilon}$ we have
\begin{eqnarray*}
\Theta_{\pi^+}(\g \cdot \gamma) &=& c_0(\pi) - \frac{\sgn_{\varepsilon}(\varpi^{-1} \mathrm{Im}_\varepsilon(
\g \cdot \gamma))}{D^G(\g \cdot \gamma)}
%&=& c_0(\pi) + f(\gamma) 
%\Theta_\pi(\g \cdot \gamma) &=& c_0(\pi) +  \sgn_\varepsilon(\varpi^{-1} \mathrm{Im}_\varepsilon}(\gamma))
%D^G(\gamma)^{-1}
\end{eqnarray*}
by Theorem 14.20 in \cite{ADSS}.   Using \ref{SGN}, we obtain
\[
\Theta_{\pi^+}(\g \cdot \gamma) = c_0(\pi) + f(\gamma)
\]
With $\gamma \in T^\varepsilon$ we also have
\begin{eqnarray*}
%Again using \cite[Theorem 14.20]{ADSS} this time with $\theta = \varepsilon, \eta = \varpi$ we have
\Theta_{\pi^-}(\gamma)
&=& \Theta_{\pi^+}(\g \cdot \gamma)
\end{eqnarray*}
and so
\[
\Theta_{\pi^-}(\gamma) = -1 + f(\gamma).
\]

The virtual character is given by 
\begin{eqnarray*}
\Theta_{\phi, s}(\gamma) &=& \rho_0(s)\Theta_{\pi^+} (\gamma) + \rho_1(s)\Theta_{\pi^-}(\gamma) \\
&=& \Theta_{\pi^+}(\gamma) - \Theta_{\pi^-}(\gamma)\\
&=& -1 - f(\gamma) - (-1 + f(\gamma))\\
&=& -2f(\gamma) .
\end{eqnarray*}
\end{proof}


 
 
  \section{Endoscopy: the non-regular parameter}  We recall from \S4 the definition of the non-regular parameter $\phi$.   The defining equations are
  \[
  \phi(\sigma) = s_2, \quad \quad \phi(\tau) = s_1.
  \]
  
  We have     
   \begin{thm}  Let $E = F(\sqrt \varepsilon)$ and let $H$ denote the norm $1$ subgroup of $E^\times$.   Let $\psi_0$ denote the quadratic character of $H$.   The non-regular parameter  $\phi$  factrorizes as follows:
     \[
  \begin{CD}
 \phi :  \bW_F @> \phi^H >> \widehat{H}  \rtimes \bW_F    @> {}^L\eta >> \widehat{G} \times \W_F
  \end{CD}
  \]
and the stable character is given by     
     \[
     \mathscr{S}\Theta_{\phi^H} = \psi_0.
     \]
      \end{thm}
      
      \begin{proof} We consider the identity component  of the centralizer of $s_1$ in $\widehat{G}$.   This is a 
      complex torus of dimension $1$.    Following \cite[\S6.2]{T} we will set 
      $\widehat{H} = \mathrm{Cent}(s_1, \widehat{G})^0$ and consider $\mathcal{H} = \widehat{H} \cdot \phi(\bW_F)$ and the inclusion
      \[
    {}^L\eta :  \mathcal{H} \to {}^LG.
      \]    
  Note that 
  \[
  \widehat{H}\cdot \phi(\bW_F) = \widehat{H} \cdot \phi(\sigma)
  \]
  since $\phi(\tau) \in \widehat{H}$.   The maximal compact subgroup of $\PGL_2(\C)$ is $\SO_3(\R)$ and $s_1, s_2$ 
  may be viewed as rotations of order $2$ about orthogonal axes.   Then any $z \in \widehat{H}$ is a  rotation about the axis of $s_1$ and the conjugate of $z$ by the rotation $s_2$ is precisely the inverse of the rotation $z$.   
  %\[  h = \left(  \begin{array}{cc}  1 & 0\\  0 & z   \end{array}  \right)_*    \in \widehat{H}    \]    then we have
  
 % \begin{eqnarray*}  \phi(\sigma) \cdot h &=& \n   \left(\begin{array}{cc}  1 & 0\\  0 & z  \end{array}  \right)_* \n^{-1}\\  & =& \left(  \begin{array}{cc}  1 & 0\\  0 & 1/z  \end{array}  \right)_*  \end{eqnarray*}
  
 So  $\phi(\sigma)$ acts on $\widehat{H}$ as inversion and this  determines an action of $\Gal(E/F)$ by inversion on $\widehat{H}$.   
  Let $H$ be the quasi-split group
defined over $F$ that is dual to $\widehat{H}$  and whose rational structure is determined
by $\Gal(E/F) \to  \mathrm{Out}(H)$ as above.  
The  group of norm $1$ elements in $E$ admits a canonical action of $\Gal(E/F)$ by inversion and so qualifies as the endoscopic group $H$.    
% This creates the extended endoscopic triple\[\fe = (H, s_1, {}^L \eta).\]

The parameter $\phi$ factors through ${}^LH$:
    \[
  \begin{CD}
 \phi :  \bW_F @> \phi^H >> \widehat{H}  \rtimes \bW_F    @> {}^L\eta >> \widehat{G} \times \W_F.
  \end{CD}
  \]
  
  It is worth noting that the centralizer of $s_1$ in $\widehat{G}$ is precisely ${}^LH$.   
  %as the crossed product of $\widehat{H}$ by the group $\langle s_2 \rangle$:  \[  \mathrm{Cent}(s_1, \widehat{G}) =  \widehat{H} \rtimes \langle s_2 \rangle.  \]
  
  
   
The maps $\phi^H$ and ${}^L\eta$ are given explicitly by
\[
\phi^H(\sigma) = 1 \rtimes \sigma, \quad \quad \phi^H(\tau) = s_1 \rtimes \tau
\]
and
\[
{}^L\eta (1 \rtimes \sigma) = s_2  \times \sigma, \quad {}^L\eta(s_1 \rtimes \tau) = s_1 \times \tau
\]




%\[
%{}^L\eta(z \rtimes \sigma) = \widehat{\eta}(z)\n \times \sigma, \quad \quad {}^L\eta(z \rtimes \tau) = \widehat{\eta}(z) \times \tau
%\]
%for all $z \in \widehat{H}$.   Note that $z_1\n z_2 \n =  z_1z_2^{-1}$  for all $z_1, z_2$
%and so ${}^L\eta$ is a homomorphism.   To see this, note that the maximal compact subgroup of $\PGL_2(\C)$ is
%$\SO_3(\R)$ and $s, \n$ may be viewed as rotations of order $2$ about orthogonal axes.
%Then $z$ is a rotation about the axis of $s$ and the conjugate of $z$ by the rotation $\n$ is precisely the inverse of the rotation $z$ 
%Note that
%\begin{eqnarray*}
%(\widehat{\eta}(z_1)\n \times \sigma)(\widehat{\eta}(z_2) \n \times \sigma) &=& 
%\widehat{\eta}(z_1)\n\widehat{\eta}(z^{-1})\n \times 1
%\end{eqnarray*}
  We now have the extended endoscopic triple
  \[
  \mathfrak{e}_1: = ({}^LH, s_1, {}^L\eta)
  \] 



  
 %The $1$-cocycle $\alpha$ attached to $\phi_H$ has values in $\widehat{H}$ and is defined as follows:
 %  \begin{eqnarray*}  \alpha(w) &=& 1 \quad \quad \textrm{if} \quad w|_E = \sigma\\  \alpha(w) &=& s\quad \quad \textrm{if} \quad w|_{E'} = \tau
% \end{eqnarray*}
  
  %We have the local Langlands correspondence for tori \cite[\S9]{Bor}:   \[    \Irr(H) \simeq \H^1_c(\bW_F, \widehat{H})   \]  In this correspondence, the quadratic character $\psi_0$ corresponds to the $1$-cocycle $\alpha$.     
  %\emph{This construction depends on choosing
% \footnote{The choice is not canonical, contrary to assertions in \cite[p.102--103]{Cun}}   
 % an ordered pair  of quadratic extensions $E/F$ and $E'/F$; the first is used to define
 % $T$ and the $L$-embedding ${}^L\eta$, the second to supply the  character $\psi_0$.}
 
% We therefore have \[ \Theta_{\phi_H} = \psi_0 \] and, for the stable character,  \[ S\Theta_{\phi_H} = \psi_0 \]   \end{proof} 
  
  
 % \emph{Notation}.   We will write
 % \begin{eqnarray*}
 % \fe_1 &=& ({}^LH, s_1, {}^L\eta) \; \textrm{with} \; H \; \textrm{the norm} \;  1 \; \textrm{subgroup of} \;
 %  F(\sqrt \varepsilon)\\
   %\fe_2 &=& ({}^LH, s_2, {}^L\eta) \; \textrm{with} \; H \; \textrm{the norm} \;  1 \; \textrm{subgroup of} \;
  % F(\sqrt \varpi)\\
   %\fe_3 &=& ({}^LH, s_3, {}^L\eta) \; \textrm{with} \; H \; \textrm{the norm} \;  1 \; \textrm{subgroup of} \;
  % F(\sqrt \varepsilon \varpi)   
%  \end{eqnarray*} 
  
In the notation of \cite[\S4,4]{FKS} we have $(S, \theta) = (H, \psi_0)$.    Then $\phi^H$ is the  $L$-parameter of the quadratic character $\psi_0$.   Therefore, we  have
\[
\Theta_{\phi^H} = \psi_0
\]
and, for the stable character,
\begin{equation}\label{STheta}
\mathscr{S}\Theta_{\phi^H} = \psi_0.
\end{equation}
\end{proof}
%Let $H_1$ denote the endoscopic group $\rU(1)$ split over $F(\sqrt{\varepsilon})$.
%let $H_2$ denote the endoscopic group $\rU(1)$ split over  $F(\sqrt \varpi)$, 
 %and let $H_3$ denote the endoscopic group $U(1)$ split over $F(\sqrt{\varepsilon \varpi})$.
 
%We need to compute the absolute transfer factors for the extended endoscopic triples $\fe, \fe', \fe''$.    
 
  \begin{lem}\label{Delta3} Let $\gamma$ be a regular semisimple element in $T^\varepsilon$.    Then we have
  \begin{eqnarray*}
   \Delta[\fw, \fe_1](\gamma, \gamma) & =&  - f(\gamma)\\
\delta \in H_2 \implies   \Delta[\fw, \fe_2](\delta, \gamma)  & =& 0\\
\delta \in H_3 \implies \Delta[\fw, \fe_3](\delta, \gamma)   &=& 0
  \end{eqnarray*}
  with $\fe_i = \fe(s_i)$ for $i = 1.2.3$.
  
 % so that, far from the identity, we have $\Delta(\gamma^H, \gamma) = \epsilon_L(\C, \psi)$ as defined in   \cite[3.6.4]{Ta}
  %  \item If $H = H_2$  then \[\Delta(\gamma^H, \gamma) = 0.\]\item If $H= H_3$   then 
%\[ \Delta(\gamma^H, \gamma) = 0. \] \end{itemize}
 \end{lem}
 \begin{proof}   Underlying the map $\widehat{\eta} \colon \widehat{H} \to \widehat{G}$ we have  a canonical pair of isomorphisms of $F$-tori:
   \begin{eqnarray*}
  T^\varepsilon & \to & H\\
  \left(\begin{array}{cc} 
   a & b\\
 \varepsilon b & a
 \end{array}
 \right)
   & \mapsto &  a + b \sqrt{\varepsilon}
   \end{eqnarray*}
   and
    \begin{eqnarray*}
     \left(\begin{array}{cc} 
   a & b\\
 \varepsilon b & a
 \end{array}
 \right)
   & \mapsto &  a - b \sqrt{\varepsilon}.
   \end{eqnarray*}
   These isomorphisms are admissible in the sense of \cite[\S3]{Ka2}.     
   We therefore have two related pairs: 
  \begin{eqnarray}\label{RP}
    (\gamma, \gamma)\quad \textrm{and} \quad (1/\gamma, \gamma). 
  \end{eqnarray}    
  %so that, given $\gamma$, there are two related elements, namely $\gamma^H$ and 
 % $1/\gamma^H$, see \cite[\S1]{Ka2}. 
 
 We will now quote the Kaletha formula for the absolute transfer factor \cite[Example 3.6.9]{Ka4}.  In our notation, this is
 \begin{equation}\label{Deltadef}
  \epsilon(1/2, \C, \psi) \cdot \kappa_{E/F}\left(\frac{c - \overline{c}}{2\eta}\right) |c - \overline{c}|_F .
 \end{equation}
 We recall that
 \[
 D_G(\gamma) = |c - \overline{c}|_F = |b|.
 \]
 In order to match the normalization of transfer factors in the recent article \cite{AK}, we switch the position of the Weyl discriminant $|c - \overline{c}|_F$ from numerator to denominator.  
The  transfer factor defined in \cite{AK}, for which there is a clean formulation of the endoscopic character identity, is the transfer factor in \cite[p.163]{Ka1} divided by $D_G(\gamma)^2$.
 This leads to the equation
 \begin{equation}\label{Deltadef}
\Delta[\fw, \fe_1](c,c) =  \epsilon(1/2, \C, \psi) \cdot \kappa_{E/F}\left(\frac{c - \overline{c}}{2\eta}\right) \cdot |c - \overline{c}|_F^{-1}.
 \end{equation}
 It follows that
% Now the normalization of the transfer factor in \cite{Ka1}  differs from that in \cite{AK}.  This difference is explained in sections %2.2 and 2.5 in \cite{AK}.    The formula (\ref{Deltadef}) is derived in the context of \cite{Ka1}.  To convert this formula to the %definition of transfer factor in \cite{AK}, we have to divide by $D^G(\gamma)^2$, see the first equation on p.163 of \cite{Ka1}.   
%We have \[D^G(\gamma) = |c - \overline{c}|_F \]and this leads to the following expression:
\begin{eqnarray*}
\Delta[\fw, \fe_1](c,c) 
%&=& \epsilon(1/2, \C, \psi) \cdot \kappa_{E/F}\left(\frac{c - \overline{c}}{2\eta}\right) |c - \overline{c}|_F^{-1} \\
&=& \epsilon(1/2, \C, \psi) \cdot \sgn_\varepsilon(b) \cdot |b|^{-1}.
\end{eqnarray*}

 Local class field theory allows us to replace the $\bW_F$-module $\C$ with the sign character 
 $ \sgn_\varepsilon$.  Denote this character by $\chi$.   We infer that
 \[
 \epsilon(1/2, \C, \psi) = \epsilon(\chi, 1/2, \psi)
 \]
 the Tate local constant of $\chi$ as in \cite[\S23.4]{BH}.   Note that ``$\Lambda$ has depth zero" in \cite{ADSS} is the same as 
 ``$\psi$ has level one"  in \cite{BH}.   We now apply the Proposition in \cite[\S23.5]{BH} and we obtain
 \begin{eqnarray*}
 \epsilon(1/2, \C, \psi) &=& \epsilon(\chi, 1/2, \psi)\\
 &=& \chi(\varpi)^{-1}\\
 &=& -1.
 \end{eqnarray*}
Putting all this together, we conclude that 
\begin{equation}\label{Delta7}
\Delta[\fw, \fe_1](\gamma , \gamma) = - f(\gamma).
\end{equation}


According to (\ref{fGalois}), $f$   is invariant under the action of $\Gal(E/F)$, from which we infer that 
\begin{equation*}\label{DeltaFormula}
\Delta[\fw, \fe_1](1/\gamma, \gamma) = - f(\gamma).
\end{equation*}

Set $\widehat{H_2} = \Cent(s_2, \widehat{G})^0$.   Then 
 \[
  \widehat{H_2}\cdot \phi(\bW_F) = \widehat{H_2} \cdot \phi(\tau)
  \]
 and $\phi(\tau)$ acts on $\widehat{H_2}$ by inversion.    Therefore,  $H_2$ is secured as the group of norm $1$ elements in $E_2 = F(\sqrt \varpi)$.   
 We have the extended endoscopic triple
 \[
 \fe_2 = (s_2, H_2, {}^L\eta_2).
 \]

  

Set $\widehat{H_3} = \Cent(s_3, \widehat{G})^0$.   Then 
 \[
  \widehat{H_3}\cdot \phi(\bW_F) = \widehat{H_3} \cdot \phi(\sigma) \phi(\tau)
  \]
 and $\phi(\sigma)$, $\phi(\tau)$ separately act on $\widehat{H_3}$ by inversion.    Therefore,  $H_3$ is secured as the group of norm $1$ elements in $E_3 = F(\sqrt {\varepsilon \varpi})$.   
 We  have the extended endoscopic triple
 \[
 \fe_3 = (s_3, \H_3, {}^L\eta_3).
 \]


At this point, write $H = H_1$ and $E = E_1$.   Note that $H_1$ (resp. $H_2, H_3$) is the endoscopic group $U(1)$ split over $E_1$ (resp. $E_2, E_3$), as in \cite[\S11.1.5]{CFM+}.



There is no admissible isomorphism from $T^\varepsilon$ to $H_2$.   If $\delta \in H_2$ and $\gamma \in T^\varepsilon$ then $\gamma $ and $\delta$ are not related and so 
\[
\Delta[\fw, \fe_2](\delta, \gamma) = 0.   
\]
Similarly, there is no admissible isomorphism from $T^\varepsilon$ to $H_3$.   If $\delta \in H_3$ and $\gamma \in T^\varepsilon$ then $\gamma$ and $\delta$ are not related and so 
\[
\Delta[\fw, \fe_3](\delta, \gamma) = 0.   
\]
\end{proof}

\begin{thm}\label{abc}  Let $\phi$ be the non-regular supercuspidal Langlands parameter for $\SL_2$.   If $\gamma \in T^\varepsilon$ then  we have  
\begin{eqnarray*}
\mathscr{E}(s_1:\gamma) &=& -2f(\gamma) \psi_0(\gamma)\\
\mathscr{E}(s_2:\gamma) &=& 0\\
\mathscr{E}(s_3:\gamma) &=& 0.
 \end{eqnarray*}
 \end{thm}
 \begin{proof}   We have
   \begin{eqnarray*}
   \mathscr{E}(s_1: \gamma) &=& \sum_{\delta \in H(F)} \Delta[\fw, \fe_1](\delta, \gamma) \, \mathscr{S}\Theta_{\phi^H} (\delta)\\
    &=& \sum_{\delta \in H(F)} \Delta[\fw, \fe_1](\delta, \gamma)
 \psi_0(\delta) \quad \textrm{by} \quad (\ref{STheta})\\ 
 &=&  \Delta[\fw, \fe_1](\gamma, \gamma)  \psi_0(\delta) 
  + \Delta([\fw, \fe_1](1/ \gamma, \gamma)) \psi_0( 1/\gamma)\\
  &=& - f(\gamma) \psi_0(\gamma) - f(\gamma) \psi_0(1/\gamma) \quad \textrm{by Lemma} \:\, \ref{Delta3}\\
  &=& -2 f(\gamma) \psi_0(\gamma) \quad \textrm{since} \quad \psi_0 = \psi_0^{-1}.
    \end{eqnarray*} 
    
    We also have
     \begin{eqnarray*}
   \mathscr{E}(s_2: \gamma) &=& \sum_{\delta \in H(F)} \Delta[\fw, \fe_2](\delta, \gamma) \, \mathscr{S}\Theta_{\phi^H} (\delta)\\
    &=& 0 \quad \textrm{by Lemma} \:\, \ref{Delta3} 
        \end{eqnarray*} 
 and       
          \begin{eqnarray*}
   \mathscr{E}(s_3: \gamma) &=& \sum_{\delta \in H(F)} \Delta[\fw, \fe_3](\delta, \gamma) \, \mathscr{S}\Theta_{\phi^H} (\delta)\\
    &=& 0 \quad \textrm{by Lemma} \:\, \ref{Delta3}
            \end{eqnarray*} 
  \end{proof} 
 

 Far from the identity, the right-hand-side of Theorem \ref{abc} is  
\[
 - 2\psi_0(\gamma) 
\]
since $f(\gamma) = 1$ by Lemma \ref{farfrom}.   Near the identity, the right-hand-side of Theorem \ref{abc} is 
\[
 - 2f(\gamma)
\]
since $\gamma$ and $1/\gamma$ are both near the identity and $\psi_0$ has depth zero.   This mirrors the calculations in \cite{ADSS}: in that paper, the character formulas are  done separately, according as $\gamma$ is far from the identity or near the identity.   
%Note that Theorem \ref{A} fits into the framework of \cite{Ka2} since $\phi$ is a regular parameter. 


  
\section{Endoscopy: regular parameters}  
%Let $\phi$ be a regular parameter and let $\psi$ be the correspoding regular character of the elliptic torus $T^\varepsilon$.   
 We recall the elements $s_1$, $s_2$ and $s_3$ in $\PGL_2(\C)$ defined in \eqref{eqn:si}.
Recall the unramified quadratic extension  $E/F$ and  
let $\sigma$ be the generator of $\Gal(E/F)$.  Let $\xi$ be a character of $E^\times$ such that 
\[
\xi \neq \xi^\sigma \quad \textrm{and} \quad \psi = \xi|_{E^1}.
\] 
Then $(E/F, \xi)$ is an admissible pair in the sense of \cite[\S18]{BH}.   Let 
\[
\rho_\xi = \Ind_E^F \, \xi
\]
be the induced $\W_F$-module.   This is a $2$-dimensional imprimitive $\W_F$-module.    We describe this module explicitly.
%Let $\pi_\xi$, via the Langlands correspondence, be the corresponding irreducible supercuspidal 
%representation of $\GL_2(F)$. 

We note that $\W_F/\W_E \simeq \Gal(E/F)$.   Coset representatives are $\{1, \sigma\}$.   

As a vector space, the induced module is $\C \oplus \C$ with module structure as follows:
\begin{itemize}
\item $w \in \W_E$ acts on the first summand $\C$ as $\xi(w)$ 
\item $w \in \W_E$ acts on the second summand $C$ as $\xi^\sigma(w)$
\item $\sigma$ sends $(z_1, z_2)$ to $(z_2, z_1)$.
\end{itemize}
This $\W_F$-module determines a homomorphism
\begin{eqnarray*}
\rho_\xi \colon\W_F &\to& \GL_2(\C)\\
w &\mapsto& \left(\begin{array}{cc} \xi(w) & 0 \\
0 & \xi^\sigma(w)
\end{array}
\right)\quad \textrm{if} \quad w \in \W_E\\
\sigma &\mapsto& \left(\begin{array}{cc}
0 & 1\\
1 & 0
\end{array}
\right).
\end{eqnarray*}

%\begin{eqnarray}\label{CD1}\rho_\xi : \W_E \rtimes \Gal(E/F) & \to & {}^LG \\w \rtimes \sigma  & \mapsto & \left(\begin{array}{cc} \xi(w) & 0\\0 & \xi^\sigma(w)\end{array}\right) \rtimes \sigma%(1,\sigma)  & \mapsto & s_2\end{eqnarray}where $\sigma$ acts as $\Ad\, s_2$.    
Then  $\phi$  is  the composition
\begin{equation}\label{CD2}
\begin{CD}
\phi \colon W_F @>\rho_\xi >> \GL_2(\C) @ >>> \PGL_2(\C).
\end{CD}
\end{equation}  
The centralizer of $\im \, \phi$ in $\PGL_2(\C)$ is $\{1, s_1\}$ and so we have
\[
S_\phi = \Z/2\Z.
\]

%Explicitly, $\rho_\xi$ is given as follows:
%Let $\varphi$ be an irreducible two-dimensional representation of $\W_F$ which lifts $\phi$.   According to \cite[\S3.3]{AMPS}, there exists a separable quadratic %extension $E/F$ and a character $\xi$ of $E^\times$ such that $\varphi \simeq \Ind_{E/F} \, \xi$.   Since $\pi$ has depth-zero, the extension $E/F$ is unramified (see %\cite[\S19.1]{BH}).  By the irreducibility of $\varphi$, we have $\xi^\sigma \neq \xi$ where $\sigma$ is the generator of $\Gal(E/F)$.   



We consider the representations $\pi:=\pi^+=\pi(\bbT^\epsilon,\psi)$ and $\pi^-$ defined in \eqref{eqn:reg}.
%Let $\phi$ denote an $L$-parameter for $\pi^+$.   Then $S_\phi = \Z/2\Z$ and 

\begin{lem}We have
\[
\Pi_\phi (\SL_2(F))= \{\pi, \pi^-\}.
\]
\end{lem}
\begin{proof}  Given the admissible pair $(E/F, \xi)$ the construction of the irreducible supercuspidal representation $\pi_\xi$ is described 
in \cite[19.2]{BH}.   Then we have the Langlands correspondence for $\GL_2(F)$ in \cite[34.4]{BH}:
\[
\rho_\xi \mapsto \pi_{\Delta_{\xi}\xi}
\]
We note that, since $E/F$ is an unramified quadratic extension, the auxiliary character $\Delta_\xi$ is unramified of order $2$ by the Definition in 
\cite[34.4]{BH}. 
Also $\kappa_{E/F}$ is unramified of order $2$ by the Proposition in \cite[34.3]{BH}.  Therefore we have  $\Delta_\xi = \kappa_{E/F}$ and twisting by $\Delta_\xi$ has no effect on $\rho_\xi$ by the Lemma in \cite[34.1]{BH}.  So the Langlands correspondence is given by
\[
\rho_\xi \mapsto \pi_\xi.
\]
Now we refer to the definition of $\pi^+$ and $\pi^-$ in (\ref{eqn:reg}).  
\end{proof}  
Let  $\rho_0$ denote the trivial character of $S_\phi$ and let $\rho_1$ denote the non-trivial character of $S_\phi$.   
We will assume that $\pi^+$ is $\fw$-generic.   Then we have a bijection
\[
\iota \colon\Pi_\phi(\SL_2(F)) \to \Irr(S_\phi), \quad \quad  \iota(\pi^+) = \rho_0, \quad \iota(\pi^-) = \rho_1.
\]


 Let $s = s_1$.   Let $\widehat{H}$ denote the identity component of the centralizer of $s$ is $\PGL_2(\C)$.   
  Define
  \[
  \mathcal{H}: = \widehat{H} \cdot \phi(W_F)
  \]
  and let ${}^L\eta : \mathcal{H} \to \PGL_2(\C)$ denote inclusion.   We recall that $E$ is the unramified quadratic extension of $F$.    Noting that $\phi(\bW_E) \subset \widehat{H}$, we have
  \[
  \mathcal{H} = \widehat{H} \cdot \phi(\sigma).
  \]
  
 
 
  This determines an action of $\Gamma_{E/F}$ by conjugation on $\widehat{H}$.    This creates the extended endoscopic triple
\[
\fe = (s, H, {}^L\eta).
\]

We note that this extended endoscopic triple is \emph{identical} to the one constructed for the non-regular parameter. In the two cases (regular parameter, non-regular parameter) the endoscopic data are the same, but the little parameters $\phi^H$ are different.   


 
 \begin{thm}\label{ENDO}  Let $\gamma$ be a regular semisimple element in $T^\varepsilon$.      We have 
 %Let $\gamma$ be a regular semisimple element in the elliptic torus $T^\varepsilon$ which is \textbf{far from the identity}.   Let $s$ denote the non-trivial element in %$S_\phi$.   We have  \[ \Theta^s_\phi(\gamma) = -\psi(\gamma) - \psi(\gamma^{-1})  \]  and 
  \[
 \sum_{\delta \in H(F)}  \Delta[\fw, \fe](\delta, \gamma) \mathscr{S}\Theta_{\phi^H} (\delta) = 
 - f(\gamma)( \psi(\gamma) + \psi(1/ \gamma)) .
  \]
    \end{thm}
  \begin{proof} 
We have
% The non-trivial element $\sigma$ in $\Gal(E/F$ acts as inversion on $\widehat{H}$.  Define ${}^L\eta$ as follows:
  %  \begin{eqnarray*}
% {}^L\eta \colon  \widehat{H} \rtimes \Gal(E/F)  & \to & \widehat{G} \times \Gal(E/F)\\
  %(w) &=& (\xi(w),1) \quad \textrm{if}\quad w \in \W_E\\
  %(\sigma) &=& (1,\sigma)\\
 % {}^L\eta(A \rtimes 1) &=& A \times 1\\   {}^L\eta(1 \rtimes \sigma) &=& \left(\begin{array}{cc}    0 & 1\\      1 & 0\\  
   %  \end{array}\right) \times \sigma.  \end{eqnarray*}  This creates an extended endoscopic triple  \[
 %\fe_0 : =  (H, s, {}^L\eta).  \]  Then $\phi$ factorizes as follows:
\[
  \begin{CD}
 \phi \colon  \bW_F @> \phi^H >> \mathcal{H}    @>  {}^L\eta >> \PGL_2(\C)
  \end{CD}
  \]
  where the little map   $\phi^H$ is the $L$-parameter of $\psi$ as in \cite[p.1148]{Ka2}.    
    Since $\phi^H$ is the $L$-parameter of $\psi$ we have 
  $\Theta_{\phi^H} = \psi$ and therefore 
   \begin{eqnarray}\label{SS}
    \mathscr{S}\Theta_{\phi^H} = \psi.
  \end{eqnarray}

 
   The transfer factor   is given by the formula (\ref{DeltaFormula}), namely
   %   We quote the formula (\ref{DeltaFormula}) for the absolute transfer factor:
   \begin{equation}\label{DeltaFormula2}
   \Delta[\fw, \fe](\gamma,\gamma ) = \Delta[\fw, \fe](1/\gamma, \gamma) = - f(\gamma).
   \end{equation}
   

  
  Finally, we have  
   \begin{eqnarray*}
 \sum_{\delta \in H(F)} \Delta[\fw, \fe](\delta, \gamma) \, \mathscr{S}\Theta_{\phi^H} (\delta) &=& \sum_{\delta \in H(F)} \Delta[\fw, \fe](\delta, \gamma)
 \psi(\delta) \quad \textrm{by} \quad (\ref{SS})\\
 &=& \Delta[\fw, \fe](\gamma, \gamma)\psi(\gamma) 
  + \Delta[\fw, \fe](1/ \gamma, \gamma)) \psi(1/ \gamma)
 \quad \textrm{by} \quad (\ref{RP})\\
  &=& - f(\gamma)(\psi(\gamma^H)  + \psi(1/ \gamma^H)) \quad \textrm{by} \quad (\ref{DeltaFormula2})
  \end{eqnarray*}  
   as required. 
  \end{proof}
  
%  We note that the endoscopic character identity    \[\Theta_{\phi,s} (\gamma)  =  \sum_{\delta \in H(F)}\Delta[\fw, \fe](\delta, \gamma) S\Theta_{\phi^H} (\delta)  \] is essentially the character identity 2.5.2 in \cite{Ka4}.  

  

  

\section{Stability across the inner forms of $\SL_2$} 
Let $\phi$ denote the non-regular parameter of $\SL_2$.   
The finite group $S_\phi$ admits a pull-back from the adjoint group $\PGL_2(\C)$ to its simply-connected 
cover $\SL_2(\C)$.   The pull-back group $\widetilde{S}_\phi$ is the quaternion group $Q$, which admits four one-dimensional representations 
$\rho_1, \rho_2, \rho_3, \rho_4$ and one two-dimensional representation $\rho_5$.   We have a bijective map
\[
\Irr(\widetilde{S}_\phi) \simeq \{\pi_1, \pi_2, \pi_3, \pi_4, \pi_5\}.
\]
The $L$-packet $\Pi_\phi$ is therefore part of a \emph{compound packet}, comprising five representations,
where $\pi_5$ is a supercuspidal representation of the inner form $\SL_1(D)$.  Let $\Theta_5$ denote the Harish-Chandra character of $\pi_5$.

\begin{thm}\label{Theta_5} The character of the supercuspidal representation $\pi_5$  is given by
\[
\Theta_5(\delta) = \psi_0(\gamma)\]
with $\delta \in \SL_1(D)$ and $\gamma \in \SL_2(F)$ whenever $\gamma, \delta$ are related.  
\end{thm} 
\begin{proof}  We will consider $\GL_2(F)$.   Now $\GL_2(F)$ admits one inner twist, namely $\GL_1(D)$ where the division algebra $D$ has the Brauer invariant
\[
\mathrm{inv}_F (D) = 1/2 \in \Q/ \Z.
\]
Let $\pi_0$ be the depth-zero supercuspidal representation of $\GL_2(F)$ for which
 \[
 \pi_0 |_{\SL_2(F)} = \pi_1 \oplus \pi_2 \oplus \pi_3 \oplus \pi_4.
 \]
 Let $\JL$ denote the Jacquet-Langlands correspondence and let 
$\pi := \mathrm{JL}^{-1} \pi_0$.   The local Langlands correspondence for $\GL_1(D)$ is determined by the composition
 \[
 \begin{CD}
 \Irr^2(\GL_1(D)) @> \JL >> \Irr^2(\GL_2(F)) @> \LLC>> \Phi(\GL_2) .
 \end{CD}
 \]
The two elements in the set
\[
 \{\pi,  \JL(\pi) \}
\] 
therefore admit the same $L$-parameter $\phi$, and comprise a \emph{compound packet}.    The corresponding
character identity is
\begin{eqnarray}\label{chid}
\Theta_\pi (\delta)= - \Theta_{\JL(\pi)}(\gamma)
\end{eqnarray}
with $\delta \in \GL_1(D)$ and $\gamma \in \GL_2(F)$ whenever $\gamma, \delta$ are related.    This character identity remains true when we restrict to $\SL_1(D)$ and $\SL_2(F)$ respectively.   

Let $Q$ denote the classical quaternion group.   According to \cite[Lemma 12.6]{HS}, we have, as a representation space for $Q \times \SL_2(F)$, 
\[
V_{\pi(\phi)} = \bigoplus_{1 \leq j \leq 4}\rho_j \otimes \pi_j.
\]

As a representation space for $\SL_2(F)$, we therefore have
\[
V_{\pi(\phi)} = \pi_1 \oplus \pi_2 \oplus \pi_3 \oplus \pi_4.
\]

According to \cite[Lemma 12.6]{HS}, we have, as a representation space for $Q \times \SL_1(D)$, 
\[
V_{\pi(\phi)} = \rho_5 \otimes \pi_5.
\]
As a representation space for $\SL_1(D)$, we therefore have
\begin{eqnarray*}
V_{\pi(\phi)} &=& \pi_5 \oplus \pi_5.
\end{eqnarray*}
%as a representation space for the direct product $Q \times \SL_1(D)$, according to \cite[Theorem 3.3]{ABPS}.    Upon restriction to the subgroup $\{1\} \times \SL_1(D)$ we obtain\[V_{\pi(\phi)} =  \pi_5 \oplus \pi_5.\]
%as a representation space for $\SL_1(D)$.   Now $\rho_5(1)$ is the identity operator on a complex vector space of dimension $2$, and so  $\pi_5$ occurs with multiplicity $2$ in the 
%restriction from $\GL_1(D)$ to $\SL_1(D)$.   
The character identity now follows from (\ref{chid})  and we obtain
%Upon restriction, we get the character identity \[\pi_5 \sqcup \pi_1 \oplus \pi_2 \oplus \pi_3 \oplus \pi_4\]and
\begin{eqnarray}\label{1234A}
2 \Theta_5(\delta) = - (\Theta_1 + \Theta_2 + \Theta_3 + \Theta_4)(\gamma)   
\end{eqnarray}
whenever $\gamma, \delta$ are related with $\gamma \in \SL_2(F)_{\reg}, \delta \in \SL_1(D)_{\reg}$.  
%From considerations of the Jacquet-Langlands correspondence we have\[\Theta_1 + \Theta_2 + \Theta_3 + \Theta_4 = - 2\Theta_5\]
%The proof, in Theorem  \ref{FAR}, of the character formula for $\Theta_{\phi, s_1}(\gamma)$ will also  yield the following formula: 

We also have
\begin{eqnarray}\label{1234B}
\Theta_1 + \Theta_2 + \Theta_3 + \Theta_4 &=& - 2\psi_0.
\end{eqnarray}
Far from the identity, this follows from Theorem \ref{FAR-NON-REG}; near the identity, this follows from Theorem \ref{NEAR-NON-REG} since
$\psi_0$ has depth zero.



By comparing (\ref{1234A}) with (\ref{1234B})   we conclude that the character of the supercuspidal representation $\pi_5$  is given by
\[
\Theta_5(\delta) = \psi_0(\gamma)\]
with $\delta \in \SL_1(D)$ and $\gamma \in \SL_2(F)$ whenever $\gamma, \delta$ are related.   
\end{proof}

The following corollary is a special case of Theorem 4.4.4(1) in \cite{FKS}.  We give an elementary proof.
%In addition, we provide a formula for the common value of the virtual character $\mathscr{S}\Theta_{\phi, *}$.   

\begin{cor}  The virtual  character $\mathscr{S}\Theta_{\phi,*}$ is stable across the two inner forms of $\SL_2$:
\[
\mathscr{S}\Theta_{\phi, 0}(\gamma) = \mathscr{S}\Theta_{\phi, 1}(\delta)
\]
for all stably conjugate strongly regular elements $\gamma \in \SL_2(F),  \delta \in \SL_1(D)$.  
\end{cor}
\begin{proof} We have
\begin{eqnarray*}
\mathscr{S}\Theta_{\phi, 0} &=& e(\SL_2(F)) \cdot (\Theta_1 + \Theta_2 + \Theta_3 + \Theta_4)\\
&=& \Theta_1 + \Theta_2 + \Theta_3 + \Theta_4\\
&=& -2 \Theta_5\\
&=& e(\SL_1(D)) \cdot 2 \Theta_5\\
&=& \mathscr{S}\Theta_{\phi, 1}
\end{eqnarray*}
as required.   
\end{proof}

%\textbf{Data Availability Statement.}   We do not analyse or generate any datasets, because our work proceeds within a theoretical and mathematical approach.


 

\markright{\textsc{References}}
\markleft{\textsc{References}}

\begin{thebibliography}{ABPSA}
\bibitem[AK]{AK} J. Adams and T. Kaletha, ``Discrete series $L$-packets for real reductive groups", arXiv: 2409.13375, 2024.
\bibitem[ADSS]{ADSS} J.D.  Adler, S. DeBacker, P.J. Sally, and L. Spice, ``Supercuspidal characters of $\SL_2$ over a $p$-adic field," Contemp. Math. \textbf{543} (2011), 19--69.
%\bibitem[ABPS]{ABPS}  A.-M. Aubert, P. Baum, R. Plymen, and M. Solleveld,   ``The local Langlands correspondence for inner forms of $\SL_n$",   Res. Math. Sci. (2016) 3:32. 
%\bibitem[AMPS]{AMPS} A-M. Aubert, S. Mendes, R.J. Plymen, and M. Solleveld, ``On $L$-packets and depth for $\SL_2(K)$ %and its inner form", International J. Number Theory \textbf{13} (2017) 2545--2568.  
%\bibitem[Au1]{Au1} A-M. Aubert, ``Representation theory of $p$-adic reductive groups", In: \href{https://www.cambridge.org/core/books/ktheory-and-representation-theory/A4E90E08F2A71BC3BEFE6B44A90C7AFF}{\textit{K-Theory and Representation Theory}}, 
%Cambridge University Press, LMS Lecture Note Series  \textbf{494}, R. Plymen and M.H.  \c{S}eng\"un (eds.), (2024).    
\bibitem[Au2]{Au2} A.-M. Aubert, ``Bruhat-Tits buildings, representations of $p$-adic groups and Langlands correspondence", J. Algebra \textbf{656} (2024) 47--76.
%\bibitem[Bor]{Bor} A. Borel, ``Automorphic $L$-functions",  Proc. Symp. Pure Math. \textbf{33} (1979), part 2, 27--61.  
\bibitem[BH]{BH} C. J. Bushnell and G. Henniart, \textit{The local Langlands conjecture for $\GL(2)$}, Grundlehren der mathematischen Wissenschaften, Vol. \textbf{335}, Springer-Verlag, Berlin, 2006. 
\bibitem[CFM+]{CFM+}  C. L. R. Cunningham, A. Fiori,  A. Moussaoui, J. Mracek, and B. Xu, \textit{Arthur Packets for $p$-adic Groups by Way of Microlocal Vanishing Cycles of Perverse Sheaves, with Examples}, Memoirs  Amer. Math. Soc., \textbf{276}:1353 (2022).    
\bibitem[DM]{DM} F. Digne and J. Michel, Representations of finite Groups of Lie type, London Math. Soc. Student Texts \textbf{21}, Cambridge University Press, Second Edition 2020.
\bibitem[FKS]{FKS} J. Fintzen, T. Kaletha, and L. Spice, ``A twisted Yu construction, Harish-Chandra characters, and endoscopy", Duke Math. Journal \textbf{172} (2023), 2241--2301. 
%\bibitem[GGS]{GGS}  I.M. Gelfand, M.I. Graev, I.I. Pyatetskii-Shapiro,  Representation Theory and Automorphic Functions, Academic Press, 1990.  
%\bibitem[HV]{HV} Henniart, Vign\'{e}ras
\bibitem[HS]{HS} K. Hiraga and H. Saito, On $L$-Packets for Inner Forms of $\SL_n$,  AMS Memoir \textbf{1013} (2012).     
\bibitem[Ka1]{Ka1} T. Kaletha, ``Endoscopic character identities for depth-zero supercuspidal $L$-packets",  Duke Math. J. \textbf{158} (2011), 161--224.
\bibitem[Ka2]{Ka2} \bysame, ``Regular supercuspidal representations'', J. Amer. Math. Soc. \textbf{32} (2019), 1071--1170.
\bibitem[Ka3]{Ka3} \bysame, ``Supercuspidal $L$-packets",  arXiv:1912.03274, 2019. 
\bibitem[Ka4]{Ka4} \bysame, ``Lectures on the stable trace formula, with emphasis on $\SL_2$", in \textit{On the Langlands Program}, Lecture Notes Series, Institute for Mathematical Sciences, National University of Singapore, vol. \textbf{43}, 2024.
%\bibitem[Ka5]{Ka5} \bysame, ``On the local Langlands conjectures for disconnected groups", preprint 2024.   
%\bibitem[KP]{KP} T. Kaletha and G. Prasad, Bruhat-Tits Theory, Cambridge University %Press, New Math. Monographs \textbf{44} (2023). 
%\bibitem[LL]{LL} J-P. Labesse, R.P. Langlands, $L$-Indistinguishability for $\SL(2)$, ..... 1979
%\bibitem[L]{L} J-P. Labesse, Stabilisation et germes pour $\SL(2)$ en toutes caract\'{e}ristiques, arXiv: 2411.14820.
%\bibitem[L]{L} R.P. Langlands, Singularit\'es et Transfert, Ann. Math. Qu\'ebec \textbf{37} (2013), %173--253.
%\ibitem[LS]{LS} R. Langlands and D. Shelstad, ``On the definition of transfer factors", Math. Ann. \textbf{278} (1987), 219--271.
\bibitem[SS]{SS} P.J. Sally and J.A. Shalika,  ``Characters of the discrete series of representations of $\SL(2)$ over a local field", Proc. Nat. Acad. Sci. U.S.A. \textbf{61} (1968), 1231--1237.   
\bibitem[T]{T} O. Taibi, ``The local Langlands conjecture", Proc. Symp. Pure Math., to appear.  arXiv:2510.00632v1
%\bibitem[Neu]{Neu} J. Neukirch, ``Class Field Theory", Springer 1986.
%\bibitem[Ser]{Ser} J-P. Serre, Linear representations of finite groups, Springer GTM 42 (1977).  
%\bibitem[Sh]{Sh} D. Shelstad, ``Notes on $L$-indistinguishability", Proc. Symp. Pure Math. \textbf{33} (1979), part 2, 193--203.   
%\bibitem[Sp]{Sp} L. Spice,  Fourier transforms of semisimple orbital integrals on the Lie algebra of $\SL_2$, Pacific J. Math. 254 (2011) 407--448.
%\bibitem[Ta]{Ta} J. Tate, ``Number theoretic background", Proc. Symp. Pure Math. \textbf{33} (1979), part 2, 3--26.   
\end{thebibliography}
\end{document}

  \section{Virtual characters near the identity: the non-regular parameter}  
 
 The sign character $\sgn_\varepsilon \colon F^\times \to \{\pm 1\}$ is determined by the short exact sequence
  \[
  1 \to \rN_{E/F}E^\times \to F^\times \to \Gal(E/F) \to 1.
  \]
  The notation $\mathrm{Im}_\varepsilon$ has already been defined in \S3.1.        
  We recall the definition of $f(\gamma)$ from \eqref{eqn:fgamma}.
 Throughout this section, $\gamma$ is near the identity.   
    
%We will use the Cayley transform defined on an open subset of $\fg$ by\[c\colon X\mapsto \frac{1+X/2}{1-X/2}.\]We recall from \cite[Lemma~5.5(h)]{ADSS} that $|D_\fg(Y)|=|D_G(c(Y))|$ when $Y\in\fg_{0+}$.
 \begin{thm}\label{NEAR-NON-REG}  Let $\gamma$ be a regular semisimple element 
in the elliptic torus $T^\varepsilon$ which is near the identity.  Then we  have 
 \begin{eqnarray*}
 \Theta_{\phi, 1}(\gamma) &=& - 2\\
 \Theta_{\phi, s_1}(\gamma) &=& -2f(\gamma).
 \end{eqnarray*}
 \end{thm}
 \begin{proof}
 Since $\gamma \in G^{\rss} \cap G_{0+}$, according to Proposition 13.14 in \cite{ADSS}, we have
\begin{equation} \label{eqn:1+2}
 \Theta_1(\gamma) + \Theta_2(\gamma) = \hat{\mu} ^G_{X_1^{\varepsilon, 1}}(c^{-1}(\gamma)).
 \end{equation}
 
 Let $\mathfrak{g} = \mathfrak{sl}_2$ be the Lie algebra of $\SL_2$.   Then $D_{\mathfrak{g}} : \mathfrak{g} \to F$ is defined as follows:
 \[
 D_{\mathfrak{g}} \left(
 \begin{array}{cc}
 a & b\\
 c & -a
 \end{array}
 \right)
 = 4(a^2 + bc).
 \]
 
  
  Let $\mathfrak{t}^\varepsilon$ denote the Lie algebra of $T^\varepsilon$.  Let  $Y \in \mathfrak{t}^\varepsilon$ 
 be such that $\gamma = c(Y)$ where $c$  is the Cayley transform \cite[Lemma 5.4]{ADSS}.   Then we have  
 $|D_{\mathfrak{g}}(Y)| = |D_G(\gamma)|$ by \cite[Lemma 5.5(h)]{ADSS}. In the notation of $\ddag$ on p.65 of \cite{ADSS} we have 
 %\[ [A; B_1, B_\varepsilon, B_\varpi, B_{\varepsilon \varpi}]_{\varepsilon,0}(Y):= |\varepsilon|^{1/2}A + q^{-1}\frac{1}{D^\fg(Y)}B_{\varepsilon}(\mathrm{Im}_{\varepsilon}(Y)). \]
  \[
 [-1; 1, B_\varepsilon, 0, 0]_{\varepsilon,0}(Y)=
 -1 + q^{-1}\frac{1}{|D_\fg(Y)|^{1/2}}B_{\varepsilon}(\mathrm{Im}_{\varepsilon}(Y)).
 \]
According to $(*)$ on p.66 of \cite{ADSS}, we have \footnote{we made a (surely necessary)  correction here}
\begin{equation} \label{Be}
B_\varepsilon =
q \cdot H(\Lambda', F(\sqrt \varepsilon)) \sgn_\varepsilon\\
=  -q\cdot \sgn_\varepsilon
\end{equation}


Taking $Y = c^{-1}(\gamma)$,  we obtain
\begin{eqnarray*} 
 \hat{\mu} ^G_{X_1^{\varepsilon, 1}}(Y) &=& [-1;1,B_\varepsilon,0,0]_{\varepsilon,0}(Y)\\
 &=& -1 -
 \frac{\sgn_\varepsilon(\mathrm{Im}_\varepsilon(\gamma))}{|D_G(\gamma)|^{1/2}} \\
 &=& -1 -f(\gamma)
 \end{eqnarray*}
 by \eqref{eqn:fgamma}.
 Hence, \eqref{eqn:1+2} gives
 \begin{equation} \label{clash1}
 \Theta_1(\gamma) + \Theta_2(\gamma) = - 1 - f(\gamma).
 \end{equation}
 
We now take $Y= c^{-1}(\g \cdot \gamma)$. Since $Y$  is in the Lie algebra of $T^{\varepsilon,\varpi}$ (by \eqref{eqn:tori}), we take $\theta' = \varepsilon$, $\eta = \varpi$ in the formula in \cite[3rd line  of p.66]{ADSS}  and $B_{\varepsilon}$ as in \eqref{Be}. We obtain 
 \begin{eqnarray}\label{YY}
 [-1; 1, B_\varepsilon, 0, 0]_{\varepsilon,0}(Y)&=&
 -1 + q^{-1}\frac{1}{D^\fg(Y)}B_{\varepsilon}(\varpi^{-1}\mathrm{Im}_{\varepsilon}(Y))\\
  &=& -1+f(\gamma). \nonumber
 \end{eqnarray}
 According to $(*)$ in \cite[\S15.2]{ADSS}, we have
 \begin{equation} \label{new(*)}
 \hat{\mu}^G_{X_1^{\varepsilon, 1}}(c^{-1}(\g\cdot\gamma)) = [-1; 1, B_\varepsilon, 0, 0]_{\varepsilon,0}(Y).
 \end{equation}
 
 It follows from \eqref{YY} and \eqref{new(*)} that
 \begin{equation} \label{mumu}
 \hat{\mu}^G_{X_1^{\varepsilon, 1}}(c^{-1}(\g\cdot\gamma)) = -1 + f(\gamma).
 \end{equation}
 
 We note that $\g \cdot \gamma \in G^{\rss} \cap G_{0+}$ by (\ref{ggamma}).    Now, using \cite[Theorem 13.14]{ADSS} and \eqref{mumu}, we obtain 
\begin{eqnarray} \label{FINAL}
 \Theta_3(\gamma) + \Theta_4(\gamma) &=&\Theta_1(\g \cdot \gamma)+ \Theta_2(\g \cdot \gamma)\\
 &=& \hat{\mu} ^G_{X_1^{\varepsilon, 1}}(c^{-1}(\g\cdot\gamma))\\
 & = & -1 + f(\gamma).
 \end{eqnarray}
 Therefore, using \eqref{clash1} and \eqref{FINAL}, we have
 \begin{eqnarray*} \Theta_{\phi, s_1} &=& \Theta_1 +\Theta_2 - \Theta_3 - \Theta_4\\
 &=& - 1 - f(\gamma) - (-1 +f(\gamma))\\
 &=& -2f(\gamma)
 \end{eqnarray*}
 and
 \begin{eqnarray*} 
 \Theta_{\phi, 1} &=& \Theta_1 +\Theta_2 + \Theta_3 + \Theta_4\\
 &=& - 1 - f(\gamma) -1  +f(\gamma)\\
 &=& -2
 \end{eqnarray*}
as required.
\end{proof}




 
 
\subsection{On the discrepancy}   There are two inconsistencies.
  The first inconsistency is between Proposition 13.14 and Theorem 15.2 in \cite{ADSS}.   
  We consider the sum $\Theta_1(\gamma) + \Theta_2(\gamma)=\Theta_{\pi^+}(\gamma)+\Theta_{\pi^-}(\gamma)$ by \eqref{eqn:packet}.   
   


According to Theorem 15.2 in \cite{ADSS} with $\theta' = \varepsilon$ and $\eta = 1$  we have
 \begin{eqnarray}\label{clash2}
 \Theta_1(\gamma) + \Theta_2(\gamma) = -1
  \end{eqnarray}
So there is an obvious clash between (\ref{clash1}) and (\ref{clash2}).   
   

 
 
 The second inconsistency is between the following consequence of Theorem 15.2 in \cite{ADSS}) and the endoscopic character identity:  
   Let $\gamma$ be a regular semisimple element 
in the elliptic torus $T^\varepsilon$ which is near the identity. Then we have
% \[ \Theta_{\pi^\pm} = \frac{1}{2} \big[ -1; 1, \beta(\gamma :  \varepsilon), \pm \beta(\gamma : \varpi) , \pm \beta(\gamma : \varepsilon \varpi)\big]
% \] \end{thm} Explicitly, this is
 \begin{eqnarray*}
 \Theta_1(\gamma) &=& \frac{- f(\gamma) - 1}{2}\\
\Theta_2(\gamma) & = &  \frac{f(\gamma) - 1}{2}\\
\Theta_3(\gamma) &=& \frac{f(\gamma) - 1}{2}\\
\Theta_4(\gamma) & = & \frac{- f(\gamma) - 1}{2}.
 \end{eqnarray*} 
 and
  \begin{eqnarray*}
 \Theta_{\phi,1}(\gamma) & = & -2 \\
   \Theta_{\phi, s_1}(\gamma) &=& 0\\
 \Theta_{\phi, s_2}(\gamma) &=& 0\\
 \Theta_{\phi, s_3}(\gamma)   &=& - 2f(\gamma).
 \end{eqnarray*}
 
 \emph{This clashes with the endoscopic character identity.}
 
 
 \begin{proof}  The formula for $\Theta_1(\gamma)$ follows from   \cite[Theorem 15.2]{ADSS} with, in their notation, $\theta' = \varepsilon$ and $ \eta  = 1$.   
 We have used the fact that, in their notation,  we have $H(\Lambda', F(\sqrt \varepsilon)) = -1$.   This depends on \cite[Lemma 4.2]{ADSS} and the fact that
 $\Lambda'$ is an additive character of depth zero.   
 
 The formula for $\Theta_2(\gamma)$ follows from \emph{loc. cit.},  again with $\theta' = \varepsilon$ and $\eta = 1$.   
 
 Moving on to the character $\Theta_3$, we have $\Theta_3(\gamma)=\Theta_1(\g \cdot \gamma)$. 
 Since $\gamma\in T^\varepsilon$, we have $\g\cdot\gamma\in T^{\varepsilon,\varpi}$, by \eqref{eqn:tori}.
 Thus we get
 \begin{equation}
 \Theta_3(\gamma)= \frac{1}{2} \left\{\frac {- \sgn_\varepsilon(\varpi^{-1} \mathrm{Im}_\varepsilon(\g \cdot \gamma))}{D^G(\gamma)} - 1 \right\} 
 \end{equation}
 using Theorem 15.2 in \cite{ADSS} with $\theta' = \varepsilon$ and $\eta =  \varpi$.  
  We now apply \eqref{SGN} and infer that
  \[
 \Theta_3(\gamma) = (f(\gamma) - 1)/2.
 \]
 A parallel argument leads us to
 \[
 \Theta_4(\gamma) =  (- f(\gamma) - 1)/2.
 \]
Recall \eqref{eqn:endos}. Then the last part of the theorem now follows rather readily.   For example: 
   \begin{eqnarray*}
  \Theta_{\phi, s_3} &=& \Theta_1 - \Theta_2 - \Theta_3 + \Theta_4\\
 % &=& \Theta_1 - \Theta_2 - \Theta_2 + \Theta_1\\
  %&=& 2\Theta_1 - 2\Theta_2\\
  &=& -f(\gamma)  -f(\gamma)\\
  &=& -2f(\gamma)
  \end{eqnarray*} 
    \end{proof} 

 \subsection{A corrected version of Theorem 15.2} We float the idea that, in the statement of Theorem 15.2, \emph{the correct statement is obtained by replacing the $\pm$ inside the curly bracket by $+$}.    This gives 
 \[
 \Theta_{\pi^\pm}(\gamma) = \frac{1}{2} \left\{H(\Lambda', F_{\theta'}) \frac{\sgn_{\theta'}(\eta^{-1} \mathrm{Im}_{\theta'}(\gamma))}{|D_G(\gamma)|^{1/2}} - 1\right\}
 \]  
 
With $\gamma \in T^\varepsilon, \theta' = \varepsilon, \eta = 1$ we have
 \begin{eqnarray*}
 \Theta_1(\gamma) & = & \Theta_2(\gamma) \\
 & = & ( - f(\gamma) - 1)/2
 \end{eqnarray*}
 
 With $\gamma \in T^\varepsilon, \g \cdot \gamma \in T^{\varepsilon, \varpi}, \theta' = \varepsilon, \eta = \varpi$ we have
 \begin{eqnarray*}
 \Theta_4(\gamma) & = & \Theta_3(\gamma) \\
 &=&  \Theta_1(\g \cdot \gamma)\\ 
 &=& (f(\gamma) - 1)/2
 \end{eqnarray*}
 



